\documentclass[11pt,letterpaper]{article}
\usepackage[margin=1in,headheight=14pt]{geometry}
\usepackage{amsmath,amssymb,amsthm}
\usepackage{booktabs,array,longtable}
\usepackage{microtype}
\usepackage{xcolor}
\usepackage{xurl}
\usepackage{fancyhdr}
\usepackage[colorlinks=true,linkcolor=blue!45!black,citecolor=blue!45!black,
urlcolor=blue!45!black,pdfauthor={Research prepared for Vladimir Reshetnikov},
pdftitle={Fermat Specializations in Intuitionistic Primitive Recursive Arithmetic}]{hyperref}

\newcommand{\N}{\mathbb N}
\newcommand{\Z}{\mathbb Z}
\newcommand{\Q}{\mathbb Q}
\newcommand{\ipra}{\mathsf{iPRA}}
\newcommand{\FLT}{\operatorname{FLT}}
\newtheorem{theorem}{Theorem}[section]
\newtheorem{lemma}[theorem]{Lemma}
\newtheorem{proposition}[theorem]{Proposition}
\newtheorem{corollary}[theorem]{Corollary}
\theoremstyle{remark}
\newtheorem{remark}[theorem]{Remark}
\setlength{\emergencystretch}{2em}
\setlength{\parskip}{0.18em}
\setlength{\parindent}{1.2em}
\renewcommand{\arraystretch}{1.2}
\pagestyle{fancy}
\fancyhf{}
\fancyhead[L]{\small Fermat specializations in $\ipra$}
\fancyhead[R]{\small Exponents $3,4,5,7$ and their multiples}
\fancyfoot[C]{\thepage}
\renewcommand{\headrulewidth}{0.35pt}
\numberwithin{equation}{section}

\begin{document}
\hypersetup{pageanchor=false}
\begin{titlepage}
\centering
\vspace*{0.6in}
{\Large\scshape Constructive number theory\par}
\vspace{0.5in}
{\Huge\bfseries Fermat specializations\par}
\vspace{0.18in}
{\LARGE in intuitionistic primitive\par recursive arithmetic\par}
\vspace{0.4in}
{\Large Explicit proofs for $3,4,5,7$\par and every positive multiple\par}
\vspace{0.5in}
{\large Research prepared for Vladimir Reshetnikov\par}
\vspace{0.12in}
{\large 15 November 2026\par}
\vfill
\begin{minipage}{0.89\textwidth}
\textbf{Main result.} For every $n\geq3$ such that
$3\mid n$, $4\mid n$, $5\mid n$, or $7\mid n$, the equation
$x^n+y^n=z^n$ has no positive integral solution. The proofs use finite
descent, explicit primitive recursive constructions, and quantifier-free
induction. The covered exponents occupy $276$ of the $420$ residue
classes, a proportion $23/35$.

\medskip
\textbf{Formalization status.} The mathematical arguments and their
$\ipra$ audits are developed in full. The accompanying Lean sources give
definitions, interfaces, and exact internal proof targets; they are an
uncompiled blueprint, not a completed Lean formalization.
\end{minipage}
\vspace*{0.5in}
\end{titlepage}
\pagenumbering{roman}
\hypersetup{pageanchor=true}

\begin{abstract}
We give human-readable constructive proofs of Fermat's last theorem for
exponents $3$, $4$, $5$, and $7$, and prove uniformly that every positive
multiple of one of these exponents is covered. The quartic theorem is
strengthened to the impossibility of $x^4+y^4=z^2$. The cubic proof uses
explicit Eisenstein pairs; the quintic proof uses the golden quadratic
ring, including a bounded normalization of its infinitely many units; the
seventh-power proof uses an integer quartic descent with Lebesgue's
essential correction included. Each descent is converted to a finite
iteration or bounded Boolean induction. We specify the language, axioms,
decidable encodings, search bounds, and induction invariants needed by
intuitionistic primitive recursive arithmetic. Finally we give concrete
Lean definitions and a dependency plan targeting ProveIt's actual
\texttt{Prf PRAAx} derivations. Exact algebra checks support the
transcription; no finite experiment is presented as a proof of FLT, and
no completed Lean or internal proof certificate is claimed.
\end{abstract}

\tableofcontents
\clearpage
\pagenumbering{arabic}

\section{Results and proof architecture}

The target is Fermat's last theorem for a substantial, explicitly
described family of exponents, under a restricted constructive arithmetic.
The family proved here is
\[
 \mathcal E=\{n\geq3:3\mid n\ \vee\ 4\mid n\ \vee\ 5\mid n\ \vee\ 7\mid n\}.
\]
For every $n\in\mathcal E$ we prove the full assertion, with no upper
bound on the positive integers $x,y,z$. The exponent family is uniform:
it is one first-order theorem with variable $n$, not merely a metatheoretic
list of separate statements. Its exact closure proof is given in
Section~\ref{sec:coverage}.

\begin{center}
\begin{tabular}{>{\raggedright\arraybackslash}p{0.12\textwidth}>{\raggedright\arraybackslash}p{0.54\textwidth}>{\raggedright\arraybackslash}p{0.23\textwidth}}
\toprule
Exponent & Main construction & Descent parameter\\
\midrule
$3$ & Factor $x+y\omega$ in the explicit pair ring;
resolve all six possible units; extract three coprime cubes. &
Absolute value of the coordinate divisible by $3$.\\
$4$ & Apply primitive Pythagorean parametrization twice to
$x^4+y^4=z^2$. & The positive right-hand square root $z$.\\
$5$ & Extract fifth powers in $\Z[\varphi]$; use two parity-preserving
families and explicit unit normalization. & The positive coordinate $s$.\\
$7$ & Convert the Fermat equation to an integer quartic;
treat all four factor allocations, including the corrected branch. &
An integer $a\geq1$, decreasing by one.\\
\bottomrule
\end{tabular}
\end{center}

The classical arguments are reorganized for the weaker theory. In
particular, a sentence such as ``take a least solution'' does not serve as
the final justification. We give a total primitive recursive transition,
or a finite scan, and an explicit quantifier-free induction invariant.
The number of iterations is bounded before the iteration begins.

The results also exclude nonzero signed and rational solutions for the
covered exponents, and yield mixed-exponent corollaries. These extensions
are elementary substitutions and denominator clearing. No claim is made
about general FLT in $\ipra$, about every regular prime, or about
unprovability of the cases outside $\mathcal E$.

The mathematical sources serve as reference and independent comparison:
the direct cubic and quartic routes agree with familiar specializations
represented in Mathlib \cite{mathlib3,mathlib4}; the quintic argument is
compared with an inspected Lean source tree \cite{five-source}; and the
seventh-power argument reconstructs both Lebesgue's original paper and its
correction \cite{Lebesgue1840,Lebesgue1840Add}. All specialized lemmas
needed by the proofs are established below.

\section{The exact constructive theory and the descent principle}
\label{sec:foundations}

Throughout, $\N=\{0,1,2,\ldots\}$ and $\Z$ is a coded set of signed natural
numbers. For $n\geq 3$, write
\[
 \FLT(n):\quad
 \forall x,y,z\in\N\;\bigl(x>0\wedge y>0\wedge z>0
       \longrightarrow x^n+y^n\ne z^n\bigr).
\]
The main assertions concern unrestricted sizes of $x,y,z$. No finite search
over possible counterexamples is being substituted for this assertion.

\subsection{Language, axioms, and logical rules}

Our $\ipra$ is first-order intuitionistic primitive recursive arithmetic.
Its numerical language has $0$, successor $S$, and a function symbol for
each specified primitive recursive construction, of every finite arity.
Initial functions are zero, successor, and projections; functions are closed
under composition and primitive recursion. For a recursion symbol $R$,
with recursion argument first, the defining equations are
\[
 R(0,\vec a)=f(\vec a),\qquad
 R(Sn,\vec a)=g(n,R(n,\vec a),\vec a).
\]
The other defining equations are those of projections and composition.
The arithmetic axioms also include $Sx\ne0$, injectivity of $S$, and
induction for every quantifier-free formula $A(n,\vec a)$:
\[
 A(0,\vec a)\ \longrightarrow\
 \bigl(\forall n\,(A(n,\vec a)\longrightarrow A(Sn,\vec a))\bigr)
 \ \longrightarrow\ \forall n\,A(n,\vec a).
\]
The logical language has equality, $\bot$, $\wedge$, $\vee$, $\to$,
$\forall$, and $\exists$. Its rules are the usual intuitionistic introduction
and elimination rules, assumption, equality reflexivity and substitution,
and substitution instances of arithmetic axioms. Quantifier introduction
and existential elimination have their usual fresh-variable conditions.
Negation abbreviates implication to $\bot$. There is no unrestricted law of
excluded middle, double-negation elimination, choice principle, or
induction for arbitrary quantified formulas.

This is precisely the presentation documented by the types
\texttt{PR}, \texttt{PTm}, \texttt{PRAAx}, and \texttt{Prf} in the inspected
ProveIt sources \cite{proveit}. In particular, \texttt{PRA.qfInd} is the
appropriate formal rule for the inductions below. The repository's
\texttt{pra\_iff\_native} translates an already constructed derivation to
the native scalar calculus. It does not turn an arbitrary theorem about
Lean's natural numbers into an $\ipra$ derivation.

\subsection{Decidable notation and bounded constructions}

For readability we use subtraction of integers, division with remainder,
gcd, divisibility, order, finite lists, and bounded quantifiers. These are
definitional conveniences. Integers can be encoded by $2a$ for $a\geq0$
and $2a-1$ for $-a<0$. A pair can be encoded by
$(a+b)^2+a$; both components of a valid code are bounded by the code, so
decoding by bounded search is primitive recursive. Repeated pairing codes
fixed tuples. Length and entries of a finite list can instead be coded by
a length and a base-$B$ integer with $B$ larger than every entry.

Every bounded loop here has a numerical clock fixed before the loop begins.
The loop state is a tuple, its update is primitive recursive, and its final
state is defined by primitive recursion on that clock. If an algorithm
stops early, subsequent updates leave its state unchanged. Searches return
the first satisfactory candidate and a success bit, or a fixed default.
An existence proof is always accompanied by a bound that makes the success
bit equal to one.

A primitive recursive predicate $P$ is represented by a characteristic
function $\chi_P\in\{0,1\}$. For example, a bounded universal quantifier is
represented by the finite product of its characteristic values, and a
bounded existential quantifier by a bounded scan. Induction on a formula
containing such notation means induction on the corresponding
quantifier-free formula involving these characteristic functions. The
equivalences with the bounded notation follow from the same finite-product
or scan recursions. Thus decidable case splits used below are legitimate
intuitionistically.

\begin{remark}[A necessary distinction]
We are not claiming that every true primitive recursive identity is
provable in $\ipra$, or that an arbitrary terminating algorithm is primitive
recursive. We give concrete bounded constructions and their arithmetic
invariants. Nor do we assume a primitive recursive universal evaluator for
all primitive recursive function codes; there is no such evaluator. Each
fixed construction used in a proof is expanded into its own function
symbols and defining equations.
\end{remark}

\subsection{Finite descent, proved using only quantifier-free induction}

\begin{lemma}[Clocked descent]\label{lem:clocked}
Let $P(s)$ be a primitive recursive predicate on coded states, and let
$D(s)$ and $r(s)$ be primitive recursive functions. Suppose that, by
quantifier-free arithmetic and decidable case distinctions, one has proved
\[
 P(s)\longrightarrow
 \bigl(P(D(s))\wedge r(D(s))<r(s)\bigr).
\tag{D}
\]
Then $\ipra$ proves $\forall s\,\neg P(s)$.
\end{lemma}

\begin{proof}
Define $s_0=s$ and $s_{k+1}=D(s_k)$ by primitive recursion. Fix a putative
state $s$ satisfying $P(s)$, and put $R=r(s)$. Induct on $k$ with the
quantifier-free invariant
\[
 k\leq R+1\ \longrightarrow\
       \bigl(P(s_k)\wedge r(s_k)+k\leq R\bigr).
\tag{I}
\]
The case $k=0$ is immediate. If $k+1\leq R+1$, the inductive hypothesis
applies at $k$, and (D) gives
$r(s_{k+1})+1\leq r(s_k)$. Adding $k$ proves the next case of (I),
including preservation of $P$. At $k=R+1$ we obtain
$r(s_{R+1})+R+1\leq R$, impossible for natural numbers. This is a
derivation of $P(s)\to\bot$, followed by universal introduction. No
infinite sequence and no least counterexample has been postulated.
\end{proof}

One can equally use a finite search for a least counterexample inside a
box whose size is supplied by an alleged solution. The clocked version
makes the induction formula and the computation bound more visible, so it
will be used throughout.

\subsection{Elementary integer arithmetic needed by the descents}
\label{sec:integer-tools}

Here are the elementary tools, including their constructive justifications.
They will also fix what is meant by ``extract a coprime power'' later.

\begin{lemma}[Division, gcd, and prime powers]\label{lem:integer-tools}
The following operations and facts are available in $\ipra$.
\begin{enumerate}
 \item For $b>0$, division with remainder produces $a=bq+r$ with
 $0\leq r<b$. Euclid's algorithm produces a gcd and integer Bezout
 coefficients.
 \item A positive integer can be factored into primes by bounded trial
 division. If a prime $\ell$ divides $uv$, then it divides $u$ or $v$.
 Valuations of nonzero integers obey
 $v_\ell(uv)=v_\ell(u)+v_\ell(v)$.
 \item If $m\geq1$ and pairwise coprime nonzero integers $a_1,\ldots,a_t$ have product
 an $m$th power, then each $|a_i|$ is an $m$th power. When $m$ is odd,
 each signed $a_i$ is an $m$th power. If $a,b>0$ are coprime and
 $ab=2u^2$, the odd factor is a square and the even factor is twice a
 square.
\end{enumerate}
\end{lemma}

\begin{proof}
The quotient is the largest $q\leq a$ with $bq\leq a$; its existence and
remainder property follow by scanning $0,\ldots,a$. The Euclidean
remainder decreases at every active step, so $a+b+1$ steps suffice for a
gcd. At each step retain the two linear combinations of the original
inputs. The identities $\gcd(u,v)=\gcd(v,u-qv)$ and
$u-qv$ equal to the updated linear combination prove the invariant.
This also proves Bezout's identity. For negative inputs apply the positive
algorithm to absolute values and adjust the signs of its coefficients.

For factorization, scan possible divisors $d=2,\ldots,a$. At a given
$d$, divide out at most $a$ copies. The invariant states that the product
of extracted factors times the current residual is $a$, and no previously
processed $d$ divides the residual. Whenever a divisor is first extracted
it is prime: a proper prime divisor would already have been removed.
That assertion is proved using the bounded search for the least divisor
of $d$ greater than one; the least divisor is prime because a nontrivial
factor of it would be smaller. Every assertion quantifies only over
integers bounded by $a$. Thus the finite invariant can be written with
characteristic functions and proved by quantifier-free induction.

If $\ell\nmid u$, primality gives $\gcd(\ell,u)=1$ by checking the
possible divisors of $\ell$; Bezout then shows that
$\ell\mid uv$ implies $\ell\mid v$. Repeated division by $\ell$, with
clock $|u|+|v|+1$, proves the valuation identity. Applied to the
factorizations already computed, it proves uniqueness of all prime
exponents.

In a coprime product, a prime occurs in at most one factor. Its valuation
in that factor therefore equals its valuation in the product, which is
divisible by $m$. Divide every exponent by $m$ and multiply the resulting
prime powers to obtain the root. For odd $m$ restore the original sign.
For $ab=2u^2$, all odd-prime valuations are even and the valuation at $2$
is odd in exactly the even factor. This gives the last assertion.
Each extracted root is bounded by the absolute value of its input; it may
alternatively be found by testing every candidate up to that bound.
\end{proof}

All later occurrences of ``coprime'' can be checked by this gcd algorithm.
All uses of valuation concern nonzero inputs and a fixed prime, or a
bounded list of prime factors of a displayed integer. Assertions such as
``every prime divisor has this property'' can therefore be read as bounded
assertions about that list.

\begin{lemma}[Primitive reduction]\label{lem:primitive}
For $n\geq1$, a positive solution of $x^n+y^n=z^n$ gives a positive
solution with pairwise coprime coordinates. Likewise, a positive solution
of $x^4+y^4=z^2$ gives one with $\gcd(x,y)=1$.
All divisions are bounded and the relevant right-hand coordinate does not
increase.
\end{lemma}

\begin{proof}
For the first equation take $d=\gcd(x,y)$. Prime valuations show
$d\mid z$: if $\ell^e\mid d$, then $\ell^{ne}\mid z^n$, so
$\ell^e\mid z$. Divide all three coordinates by $d$ and cancel $d^n$.
The new first two coordinates are coprime. A common prime of either of
them and the third divides the remaining one by the equation, a
contradiction. For the second equation the same argument gives
$d^2\mid z$; divide $x,y,z$ by $d,d,d^2$, respectively. The gcd algorithm
and the power-divisibility facts in Lemma~\ref{lem:integer-tools} give all
these witnesses constructively.
\end{proof}


\section{Exponent four: an explicit Pythagorean descent}
\label{sec:four}

We prove the stronger assertion that
\begin{equation}\label{eq:fermat42}
 x^4+y^4=z^2
\end{equation}
has no solution in positive integers. This is Fermat's classical descent,
here written with a concrete smaller state and a finite clock.

\begin{lemma}[Primitive Pythagorean parametrization]\label{lem:pythagoras}
Suppose $A,B,C>0$, $A^2+B^2=C^2$, $\gcd(A,B)=1$, and $B$ is even.
Then $A,C$ are odd and there are coprime integers $m>n>0$ of opposite
parity such that
\[
 A=m^2-n^2,\qquad B=2mn,\qquad C=m^2+n^2.
\]
The parameters can be found by a search bounded by $C$.
\end{lemma}

\begin{proof}
An odd square is $1$ modulo $4$. Thus two odd legs are impossible, and
coprimality excludes two even legs. The designated even leg makes $A,C$
odd. Put $U=(C+A)/2$ and $V=(C-A)/2$, positive integers with
$UV=(B/2)^2$. A common divisor of $U,V$ divides $C,A$. Any prime
dividing $C,A$ divides $B^2$, hence $B$, contradicting coprimality. Thus
$U,V$ are coprime squares: $U=m^2$, $V=n^2$ by
Lemma~\ref{lem:integer-tools}. Since $U>V$, $m>n>0$; since $C$ is odd,
their parities differ. Finally $B^2=4m^2n^2$ and positivity give
$B=2mn$. Root search up to $C$ finds both parameters. This proof gives
the displayed identities, rather than invoking a classification theorem
with an unbounded choice of parameters.
\end{proof}

\begin{theorem}[Fermat's stronger quartic theorem]\label{thm:four}
The theory $\ipra$ proves
\[
 \forall x,y,z\in\N\quad x>0\wedge y>0
          \longrightarrow x^4+y^4\ne z^2.
\]
Consequently it proves $\FLT(4)$.
\end{theorem}

\begin{proof}
Start with a putative solution. Its right-hand coordinate $z$ is positive
because $x,y>0$; apply primitive reduction. Exchange
$x,y$ if necessary so that $x$ is odd and $y$ even. The parity conclusion
follows directly from fourth powers modulo $16$: two odd fourth powers
sum to $2$, which is not a square residue, and coprimality excludes two
even coordinates. Lemma~\ref{lem:pythagoras}, applied to $x^2,y^2,z$,
gives
\begin{equation}\label{eq:four-first}
 x^2=m^2-n^2,\qquad y^2=2mn,\qquad z=m^2+n^2,
\end{equation}
where $m>n>0$ are coprime and have opposite parity.
If $m$ were even and $n$ odd, then $x^2\equiv-1\pmod4$, impossible.
Hence $m$ is odd and $n$ even. The equality $y^2=2mn$ and coprimality
give
\[
 m=u^2,\qquad n=2v^2,
\]
with $u,v>0$, $u$ odd, and $\gcd(u,v)=1$. More explicitly,
$m(n/2)=(y/2)^2$, the two factors are coprime, and both are squares.
Substitute these expressions into the first equation of
\eqref{eq:four-first}:
\[
 x^2+(2v^2)^2=(u^2)^2.
\]
The two legs are coprime. Indeed, a common odd prime divisor of $x,v$
would divide $u$ by this equation, while $2\nmid x$.
Apply Lemma~\ref{lem:pythagoras} a second time. There are coprime
$r>s>0$ with
\[
 x=r^2-s^2,\qquad 2v^2=2rs,\qquad u^2=r^2+s^2.
\]
Now $rs=v^2$ and $\gcd(r,s)=1$, so $r=a^2$ and $s=b^2$ for
positive coprime integers $a,b$. The final identity becomes
\[
 a^4+b^4=u^2.
\]
Thus $(a,b,u)$ is another positive primitive solution of
\eqref{eq:fermat42}. Its last coordinate is strictly smaller:
\[
 u\leq u^4=m^2<m^2+n^2=z.
\]

This construction uses only gcd, division, and bounded root searches.
One can define its total step function on every coded triple by first
checking the state predicate, then performing the indicated searches, and
returning a default on an invalid state. On a valid state each search
succeeds by the preceding lemmas. Let the state predicate be positivity,
coprimality, and \eqref{eq:fermat42}, with rank $z$; orient the two legs
inside the step function. Lemma~\ref{lem:clocked} now yields a
contradiction after at most $z+1$ iterations. Finally a solution of
$x^4+y^4=z^4$ would give \eqref{eq:fermat42} with right-hand coordinate
$z^2$.
\end{proof}

\paragraph{Numerical bounds for formalization.}
In a valid state, $x,y\leq z$. The first Pythagorean parameters are at
most $z$; the roots $u,v$ and the second parameters $r,s,a,b$ are also at
most $z$. All searches can therefore scan the box $[0,z]^2$ or the
interval $[0,z]$. Products used in checking the identities are bounded by
a fixed polynomial in $z$, for example $8(z+1)^8$ is more than enough.
The outer descent has the fixed clock $z+1$. This makes both the state
transition and the invariant in Lemma~\ref{lem:clocked} explicit
primitive recursive constructions.


\section{The cubic exponent: an explicit Eisenstein descent}
\label{sec:three}

We prove the stronger signed statement
\[
  xyz\ne0\quad\Longrightarrow\quad x^3+y^3+z^3\ne0
  \qquad(x,y,z\in\mathbb Z).
\]
Signed integers here are finite codes for a sign and an absolute value.
In particular, neither complex numbers nor an infinite number field is
needed.  The descent decreases the absolute value of the coordinate
divisible by~$3$.  Its construction and its termination both have
primitive-recursive bounds.

\subsection{The finite arithmetic needed below}

For clarity, we record the parts of Lemma~\ref{lem:integer-tools} used
below, in their algorithmic forms.  Trial division through $2,\ldots,|m|$
produces a prime factorization of each nonzero $m$.  Euclid's algorithm
produces $\gcd(u,v)$ and Bezout coefficients.  The usual proofs use division
with remainder and the identity preserving the gcd at one Euclidean
step.  They consequently use induction only over explicitly finite
computations.  Comparing prime exponents in the resulting lists proves:

\begin{enumerate}
\item If $d^3\mid m^3$, with $d>0$, then $d\mid m$.
\item If $u,v,w$ are pairwise coprime signed integers and $uvw$ is a cube,
then each of $u,v,w$ is a signed cube.  Indeed, a prime occurs in at most
one of their absolute values, so its exponent there is a multiple of
three; attach the original sign after taking the positive cube root.
\item If $\gcd(u,v)=1$ and $u^3+v^3+w^3=0$, the three coordinates are
pairwise coprime.  A prime dividing any two divides the third.
\end{enumerate}

These facts do not require an unbounded search for a factor or a root.
Divisors of $m$ are searched below $|m|$, and a cube root of a nonzero
integer $m$ is searched in $[-|m|,|m|]$.  Later applications always supply
the relevant nonzero integers.

\subsection{The pair ring and its Euclidean algorithm}

Let $E$ be the set of pairs of signed integers, written $a+b\omega$, with
\[
 \omega^2+\omega+1=0.
\]
This is a notation for the following explicit operations:
\begin{align*}
 (a,b)+(c,d)&=(a+c,b+d),\\
 (a,b)(c,d)&=(ac-bd,\ ad+bc-bd),\\
 \overline{(a,b)}&=(a-b,-b),\\
 N(a,b)&=a^2-ab+b^2.
\end{align*}
Direct multiplication verifies
\[
 \alpha\bar\alpha=N(\alpha),\qquad
 N(\alpha\beta)=N(\alpha)N(\beta),\qquad
 2N(a,b)=a^2+b^2+(a-b)^2.
\]
Thus $N\ge0$, and $N(\alpha)=0$ exactly when $\alpha=0$.
In particular, this pair ring is an integral domain: if
$\alpha\beta=0$, multiplication of natural-number norms gives
$N(\alpha)N(\beta)=0$, and decidable arithmetic gives $\alpha=0$ or
$\beta=0$.  The last displayed identity also gives the useful bound
\begin{equation}
 |a|,\ |b|,\ |a-b|\ \le 2N(a,b).
 \label{eq:cubic-coordinate-bound}
\end{equation}
Here we used $|t|\le t^2$ for every signed integer $t$.

We next give the division algorithm.  If $\beta=c+d\omega\ne0$, set
$D=N(\beta)>0$.  In the rational pair ring the quotient $\alpha/\beta$,
for $\alpha=a+b\omega$, has coordinates
\[
 u=\frac{ac-ad+bd}{D},\qquad v=\frac{bc-ad}{D}.
\]
Choose integers $q_0,q_1$ with $|u-q_0|,|v-q_1|\le1/2$, using signed
integer division and a fixed tie convention.  Put
$q=q_0+q_1\omega$ and $r=\alpha-q\beta$.  For rational numbers
$s,t$ of absolute value at most $1/2$,
\[
 s^2-st+t^2\le\tfrac14+\tfrac14+\tfrac14=\tfrac34.
\]
Consequently
\[
 N(r)\le\tfrac34N(\beta)<N(\beta).
\]
The proof is entirely rational arithmetic; denominators can simply be
cleared.  Thus Euclid's algorithm terminates after at most
$N(\beta)+1$ divisions, including a possible initial exchange.  Repeated
substitution in the division equations gives a gcd $g$ and a Bezout
identity $g=s\alpha+t\beta$.

Here are the precise consequences of this algorithm, with their proofs.
An element is a unit exactly when its norm is~$1$: necessity follows
from multiplicativity, and sufficiency from $\alpha\bar\alpha=1$.
The six units are
\[
 \{1,-1,\omega,-\omega,\omega^2,-\omega^2\}.
\]
To check completeness of this list, write
$4N(a,b)=(2a-b)^2+3b^2$.  When $N=1$, one has
$b\in\{-1,0,1\}$, and checking these three possibilities gives exactly
the six displayed elements.

Every nonzero nonunit has a factorization into irreducibles.  To construct
one, search for a nonunit proper divisor in the box
$|a|,|b|\le2N(\alpha)$, and split when one is found.  A divisor of
$\alpha$ has norm at most $N(\alpha)$, hence lies in this box by
\eqref{eq:cubic-coordinate-bound}; divisibility is checked by the explicit
quotient formula above.  If no such divisor exists, the element is
irreducible.  A proper splitting replaces a norm by two smaller positive
norms, both at least two.  Equivalently one may repeatedly split a
factor list; after $k$ proper splittings its product has at least $k+1$
nonunit factors, so $2^{k+1}\le N(\alpha)$.  In particular a loop with
$N(\alpha)$ rounds suffices.

Every irreducible is prime.  If an irreducible $\pi$ does not divide
$\eta$, its gcd with $\eta$ is a unit.  The computed Bezout identity
can then be normalized to $u\pi+v\eta=1$.  If $\pi\mid\eta\theta$,
multiplying this identity by $\theta$ proves $\pi\mid\theta$.
Finally, factorization is unique up to order and units: compare an
irreducible in one factorization with the product in the other, use
primality to find an associate, cancel it, and repeat through the finite
factor lists.  This is the only unique-factorization result used below;
it has just been proved for our explicit pair ring.

\subsection{The prime above three}

Put $\lambda=1-\omega=(1,-1)$.  Calculation gives
\[
 N(\lambda)=3,\qquad \lambda^2=-3\omega,
 \qquad \bar\lambda=-\omega^2\lambda.
\]
The norm $3$ is an ordinary prime, so any nonunit proper divisor of
$\lambda$ would have an impossible intermediate norm.  Thus $\lambda$
is irreducible and therefore prime.  For $\alpha=a+b\omega$,
\begin{equation}
 \lambda\mid\alpha\ \Longleftrightarrow\ 3\mid a+b,
 \qquad
 \lambda^2\mid\alpha\ \Longleftrightarrow\ (3\mid a\ \text{and}\ 3\mid b).
 \label{eq:cubic-lambda}
\end{equation}
The first equivalence follows from the quotient formula
\[
 \frac{a+b\omega}{1-\omega}
 =\frac{2a-b}{3}+\frac{a+b}{3}\omega:
 \quad 2a-b=2(a+b)-3b.
\]
The second follows from $\lambda^2=-3\omega$ and the fact that
$\omega$ is a unit.  These observations also prove, without any
abstract quotient-ring construction, all ramification facts needed here.

\subsection{Normalization of a hypothetical solution}

Suppose $x,y,z$ are nonzero signed integers with
$x^3+y^3+z^3=0$.  Divide by $g=\gcd(|x|,|y|)>0$.
Since $g^3\mid z^3$, the elementary cube-divisibility fact gives
$g\mid z$.  The resulting integer triple is a solution with
$\gcd(x,y)=1$, and hence is pairwise coprime.  We keep the same letters
for this primitive triple.

The cubes modulo~$9$ are $0,1,-1$.  If none of $x,y,z$ is divisible
by~$3$, the sum of their cubes is one of $-3,-1,1,3$ modulo~$9$, never
zero.  Thus one coordinate is divisible by~$3$, and pairwise coprimality
shows that exactly one is.  Permute the coordinates to arrange
\begin{equation}
 3\mid z,\qquad 3\nmid xy,\qquad 3\mid x+y.
 \label{eq:cubic-normalization}
\end{equation}
In fact $9\mid x+y$.  Write $x+y=3k$.  Then
\[
 -z^3=x^3+y^3=(x+y)((x+y)^2-3xy)
                  =9k(3k^2-xy).
\]
The left side is divisible by~$27$, while $3k^2-xy$ is not divisible
by~$3$.  Cancellation of $9$ and primality of~$3$ give $3\mid k$.
Also $x+y\ne0$, since otherwise the original equation would give $z=0$.

\subsection{Extracting the cube in the pair ring}

Set $\alpha=x+y\omega$.  It is nonzero, and
\begin{equation}
 (x+y)(x+y\omega)(x+y\omega^2)=-z^3.
 \label{eq:cubic-factorization}
\end{equation}
By \eqref{eq:cubic-lambda} and
\eqref{eq:cubic-normalization}, $\lambda$ divides $\alpha$ exactly once:
$\lambda\mid\alpha$ but $\lambda^2\nmid\alpha$.

Any irreducible $\pi$ not associate to $\lambda$ that divides $\alpha$
divides neither of the other two factors on the left of
\eqref{eq:cubic-factorization}.  Indeed, a common divisor of $\alpha$
and $x+y$ divides their difference $-y\lambda$.  A common divisor of
$\alpha$ and $\bar\alpha=x+y\omega^2$ divides
$y(\omega-\omega^2)=y\omega\lambda$.  In either case primality and
$\pi\nmid\lambda$ give $\pi\mid y$; subtracting $y\omega$ from
$\alpha$ gives $\pi\mid x$.  Ordinary Bezout coefficients for
$\gcd(x,y)=1$ would then make $\pi$ divide~$1$, a contradiction.

Compare the finite irreducible factorizations of both sides of
\eqref{eq:cubic-factorization}.  For every irreducible other than
$\lambda$, its exponent in $\alpha$ is its entire exponent on the
left, and so is a multiple of three.  The exponent of $\lambda$ in
$\alpha$ is one.  Therefore there are a unit $\varepsilon$ and an
element $\gamma=a+b\omega$ such that
\begin{equation}
 x+y\omega=\varepsilon\lambda(a+b\omega)^3,
 \qquad \lambda\nmid(a+b\omega).
 \label{eq:cubic-extraction}
\end{equation}
This is a constructive extraction from a factor list.  Alternatively,
one can find the same witness by a bounded search: taking norms shows
$3N(\gamma)^3=N(\alpha)$, so
$|a|,|b|\le2N(\alpha)$; search this box and the six units.

\subsection{The unit calculation, including the exceptional possibilities}

Write
\[
 A=a^3-3ab^2+b^3,\qquad B=3ab(a-b).
\]
Direct multiplication gives
\[
 (a+b\omega)^3=A+B\omega,
 \qquad
 \lambda(a+b\omega)^3=(A+B)+(2B-A)\omega.
\]
The sums of the two coefficients after multiplication by the three
units $1,\omega,\omega^2$ are, respectively,
\[
 3B,\qquad 3(A-B),\qquad -3A.
\]
Multiplication by their negatives just changes the signs.  Now
$3\mid B$, while
\[
 A\equiv a+b\not\equiv0\pmod3;
\]
the final inequality is exactly $\lambda\nmid\gamma$ from
\eqref{eq:cubic-lambda}.  Since $9\mid x+y$, the choices
$\varepsilon=\pm\omega,\pm\omega^2$ are impossible.  Consequently
\begin{equation}
 \varepsilon\in\{1,-1\},\qquad
 x+y=\varepsilon\,9ab(a-b).
 \label{eq:cubic-unit}
\end{equation}
This check is essential: discarding units without this argument would
leave a gap in the descent.

\subsection{The smaller solution}

Put $M=N(a,b)>0$.  Norms in \eqref{eq:cubic-extraction} give
$N(x+y\omega)=3M^3$.  Together with
\eqref{eq:cubic-unit}, this gives
\begin{equation}
 -z^3=(x+y)N(x+y\omega)
      =\varepsilon\,27M^3ab(a-b).
 \label{eq:cubic-product}
\end{equation}
Thus $(3M)^3\mid z^3$, and hence $3M\mid z$.  Set
$h=z/(3M)$.  Equation \eqref{eq:cubic-product} now says
\[
 ab(a-b)=-\varepsilon h^3=(-\varepsilon h)^3.
\]
None of $a,b,a-b,h$ is zero, since $x+y$ and $z$ are nonzero.

The integers $a,b$ are coprime.  If a positive integer $d$ divides both,
then $d^3$ divides both coordinates of
$\varepsilon\lambda(a+b\omega)^3=x+y\omega$; hence $d^3$ divides
$\gcd(x,y)=1$, so $d=1$.  It follows that $a,b,a-b$ are pairwise
coprime.  The ordinary integer cube-splitting lemma therefore supplies
nonzero signed integers $r,s,t$ with
\[
 a=r^3,\qquad b=s^3,\qquad a-b=t^3.
\]
We have obtained another cubic solution,
\[
 s^3+t^3+(-r)^3=0.
\]
It is primitive: any prime dividing two of $r,s,t$ would divide two of
$a,b,a-b$.  By the same modulo-$9$ calculation as before, exactly one
of $s,t,-r$ is divisible by~$3$.  Permute these three entries so that
this coordinate is called $z'$.

The decrease is explicit.  Since $|u|\le|u|^3$ for nonzero integers,
\[
 |r|\le|a|,\qquad |s|\le|b|,\qquad |t|\le|a-b|.
\]
By \eqref{eq:cubic-coordinate-bound} each right side is at most $2M$.
On the other hand, $z=3Mh$ with $|h|\ge1$, so
\begin{equation}
 0<|z'|\le2M<3M\le|z|.
 \label{eq:cubic-descent}
\end{equation}
We have therefore constructed, from every normalized primitive
solution, another such solution with a smaller positive distinguished
coordinate.

\subsection{Finite termination in iPRA}

It remains to instantiate Lemma~\ref{lem:clocked} explicitly, rather
than appeal to an infinite descending sequence or to a least
counterexample.  Let $V(S)$ be the
primitive-recursive Boolean predicate saying that a coded triple $S$
has nonzero pairwise coprime signed coordinates, their cubes sum to
zero, and the third coordinate alone is divisible by~$3$.  Let
$m(S)$ be the absolute value of this third coordinate.  The algorithms
above define a total primitive-recursive function $D$: when $V(S)$ is
true, perform the bounded cube extraction and subsequent bounded
integer root searches, and return the first resulting normalized
triple in a fixed order; otherwise return a fixed dummy triple.  The
proof above establishes
\[
 V(S)=1\quad\Longrightarrow\quad
 V(D(S))=1\ \text{ and }\ 0<m(D(S))<m(S).
\]
Define $S_0=S$ and $S_{k+1}=D(S_k)$ by primitive recursion.  If
$m=m(S)$ and $V(S)=1$, quantifier-free induction on $k$, with the
guard $k\le m$, proves
\[
 k\le m\quad\Longrightarrow\quad V(S_k)=1\ \text{ and }\ m(S_k)+k\le m.
\]
At $k=m$ this gives $m(S_m)\le0$, whereas $V(S_m)=1$ gives
$m(S_m)>0$.  This is a contradiction in decidable natural-number
arithmetic.  The induction formula contains only primitive-recursive
terms and decidable relations; bounded conjunctions, if used to check
factor lists, are represented by Boolean folds.  No excluded middle
for arbitrary formulas and no unbounded minimization have been used.

For an explicit size audit, if both Euclidean inputs have norm at most
$H\ge1$, their coordinates are bounded by $2H$.  In the quotient
formula, the numerators are bounded in absolute value by $12H^2$,
so rounded quotient coordinates are bounded by $Q=12H^2+1$.
There are at most $H+2$ steps.  Bezout coefficients are therefore
bounded, for example, by
\[
 (3Q+2)^{H+3}.
\]
Indeed a coefficient update is an old coefficient minus a product of
a quotient and another coefficient, and pair multiplication increases
the coordinate bound by at most a factor of $3Q$.
Factorization uses at most $N(\alpha)$ factors and rounds; each
candidate factor has coordinates bounded by $2N(\alpha)$.
The cube-extraction box has the same bound.  These coarse bounds are
all primitive recursive and suffice for coding every invariant by
bounded searches and finite lists.  The descent itself takes only
$m(S)$ steps; by \eqref{eq:cubic-descent}, after its first step all
three coordinate absolute values are below the original $m(S)$.

We have proved the signed cubic statement, and therefore
\[
 \forall x,y,z\in\mathbb N_{>0},\qquad x^3+y^3\ne z^3.
\]
Replacing $x,y,z$ by $x^k,y^k,z^k$ gives FLT for every positive
multiple $3k$ of~$3$.  This final substitution is uniform in~$k$,
since exponentiation itself is primitive recursive.

\subsection{Lean obligations for this proof}

A direct formalization should use a pair type with signed-integer
coordinates, rather than import a general number-field construction.
The following are the exact successive obligations.

\begin{enumerate}
\item Define pair addition, multiplication, conjugation and norm by the
four displayed formulas.  Prove their ring identities and norm
multiplicativity by polynomial normalization, and prove
\eqref{eq:cubic-coordinate-bound} by integer inequalities.
\item Implement signed nearest-integer division.  Prove the Euclidean
remainder bound after clearing denominators.  Define the gcd computation
with explicit norm fuel and prove its Bezout invariant.
\item Implement bounded irreducible factorization.  Prove primality by
the Bezout identity and uniqueness by finite cancellation.  Formalize
the six-unit table and \eqref{eq:cubic-lambda} as coordinate lemmas.
\item Define the Boolean predicate $V$, the bounded witness extractor
in \eqref{eq:cubic-extraction}, and the descent function $D$.  The
core theorem is the concrete implication
$V(S)=1\Rightarrow V(D(S))=1\land m(D(S))<m(S)$.
The unit table and polynomial identities above give its arithmetic
part; integer coprime cube extraction gives its final three roots.
\item Prove the guarded finite-iteration invariant and derive FLT for
$3$, then for $3k$ by the exponent law $(x^k)^3=x^{3k}$.
\item Separately translate these bounded programs, their defining
equations and their quantifier-free induction instances into the
chosen object-language presentation of iPRA.  A Lean proof about
ordinary integers alone does not perform this last step automatically.
\end{enumerate}

The existing Mathlib declaration
\texttt{fermatLastTheoremThree} in
\nolinkurl{Mathlib.NumberTheory.FLT.Three} is a useful independent
semantic comparison.  Its cyclotomic-field imports are not used as an
unexamined certificate of iPRA-provability in the argument above.

\section{The quintic specialization: a completely explicit Dirichlet descent}
\label{sec:five}

This section proves the exponent-five specialization, including both parity
cases.  Its arithmetic takes place in the fixed quadratic ring
\[
 R=\mathbb Z[\varphi],\qquad
 \varphi^2=\varphi+1,\qquad
 \varphi=\frac{1+\sqrt5}{2}.
\]
The square-root notation is an abbreviation for integer-pair calculations,
not an assumption that iPRA contains real numbers.  The proof below replaces
the usual appeal to a least counterexample by a bounded iteration.

\subsection{Arithmetic in the golden quadratic ring}

Represent \(a+b\varphi\) by the integer pair \((a,b)\).  Addition is
coordinatewise, and multiplication, conjugation, and norm are
\[
 (a,b)(c,d)=(ac+bd,\,ad+bc+bd),\qquad
 \overline{(a,b)}=(a+b,-b),\qquad
 N(a,b)=a^2+ab-b^2.
\]
Direct expansion proves
\(\alpha\bar\alpha=N(\alpha)\) and
\(N(\alpha\beta)=N(\alpha)N(\beta)\).
The norm of a nonzero pair is nonzero: otherwise
\((2a+b)^2=5b^2\), and cancelling the gcd of \(2a+b\) and \(b\)
contradicts the fact that a prime dividing a square divides its base.
All this uses only integer arithmetic.

For a convenient presentation of inequalities we use the ordered quadratic
field \(\mathbb Q(\sqrt5)\).  It too is an arithmetic coding device:
the sign of \(A+B\sqrt5\), after clearing positive denominators, is decided
by the signs of \(A,B\), and, when these have opposite signs, by comparing
\(A^2\) with \(5B^2\).  Equality in the latter comparison is impossible
unless \(A=B=0\).  The ordered-field laws follow by multiplying out these
finite sign cases.  In particular, no principle about arbitrary real
numbers, completeness, or choice is used.  The elementary inequalities
\[
 2<\sqrt5<3,\qquad \frac32<\varphi<2,\qquad
 \varphi^{-1}=\varphi-1
\]
are verified by squaring positive rational numbers.

\paragraph{Euclidean division.}
Write \(\alpha/\beta=x+y\varphi\) with rational \(x,y\) and
\(\beta\ne0\).  Choose nearest integers \(m,n\) by quotient and remainder,
breaking ties in a fixed direction.  Put
\(r=x-m\), \(s=y-n\).  Since \(|r|,|s|\le1/2\),
\[
 |N(r+s\varphi)|
 =|r^2+rs-s^2|\le r^2+|rs|+s^2\le\frac34.
\]
Consequently
\[
 \rho=\alpha-(m+n\varphi)\beta
 \quad\Longrightarrow\quad
 |N(\rho)|\le\frac34|N(\beta)|<|N(\beta)|.
\]
Thus the ordinary Euclidean algorithm works with the nonnegative integer
measure \(|N|\).  If its initial second argument has norm magnitude \(M\),
there are at most \(M+1\) divisions.  Back substitution gives a Bezout
identity for its gcd.

Here are the elementary consequences, with their constructive content
made explicit.  An element is a unit exactly when its norm is \(1\) or
\(-1\), since then \(\bar\alpha/N(\alpha)\) is its inverse.  An irreducible
element \(\pi\) is prime: if \(\pi\nmid a\), its gcd with \(a\) is a unit,
so a Bezout identity \(u\pi+va=1\), multiplied by \(b\), proves
\(\pi\mid b\) whenever \(\pi\mid ab\).  Factorization terminates because
a factorization into two nonunits splits the positive integer \(|N|\)
into two factors each at least \(2\).  Uniqueness follows by repeatedly
using this prime-divisor property and cancelling.

It is useful to explain why this argument need not conceal unbounded
search.  Let \(\alpha=a+b\varphi\ne0\) and
\(H=|a|+|b|+1\).  In the positive embedding, both
\(|\alpha|\) and \(1/|\alpha|\) are less than \(2H\): the first
bound is immediate, and the second follows from
\(|N(\alpha)|\ge1\) and \(|\bar\alpha|<2H\).
The inequality
\(\varphi^j\ge1+j/2\), proved by ordinary induction, shows that a
signed exponent \(k\) with \(|k|\le4H\) normalizes a signed associate
\[
 1\le v=\pm\alpha\varphi^k<\varphi.
\]
Explicitly, the endpoints \(\varphi^{-4H}|\alpha|<1\) and
\(\varphi^{4H}|\alpha|>1\) bracket the crossing of \(1\).
Scanning consecutive exponents and taking the first value at least \(1\)
gives the required strict upper bound \(\varphi\).
If \(M=|N(\alpha)|\), then \(|\bar v|\le M\).  Solving for the two
integer coordinates of \(v\) bounds both by \(3M+6\).
Therefore, if \(\alpha\) is reducible, a nonunit proper divisor can be
found among the finitely many pairs in the box
\([-(3M+6),3M+6]^2\), by checking divisibility and
\(2\le |N(\delta)|<M\).
Each quotient is computed by the conjugate formula and a test that its
two rational coordinates are integers.  This supplies a bounded
factorization procedure; at most \(M\) nonunit factors are needed.
One can normalize factors as they are produced.  Thus the preceding
existence and uniqueness arguments have primitive-recursive
implementations, with quantifier-free invariants on finite lists.

\paragraph{All units.}
Every unit is \(\pm\varphi^j\), \(j\in\mathbb Z\).
For completeness, normalize a unit as above to \(1\le v<\varphi\),
and write \(v=a+b\varphi\).  If \(N(v)=1\), then
\[
 0\le b=\frac{v-v^{-1}}{\sqrt5}
 <\frac{\varphi-\varphi^{-1}}{\sqrt5}
 =\frac1{\sqrt5}<1.
\]
Thus \(b=0\), and \(v=1\).  If \(N(v)=-1\), then
\[
 0<b=\frac{v+v^{-1}}{\sqrt5}<1,
\]
which is impossible.  For the last strict inequality use
\(\varphi+\varphi^{-1}=\sqrt5\) and the identity
\[
 (\varphi+\varphi^{-1})-(v+v^{-1})
   =(\varphi-v)\left(1-\frac1{\varphi v}\right)>0.
\]
Undoing the normalization proves the assertion.
In particular, every unit is \(u^5\varphi^r\) for a unit \(u\)
and a uniquely selected \(r\in\{0,1,2,3,4\}\);
the sign is absorbed in a fifth power and the exponent is divided by \(5\).

\subsection{A fifth-power extraction lemma}

We need to remove the unit in a coprime factorization.  Reduction modulo
\(5\) is especially simple:
\[
 \varphi^5=3+5\varphi,\qquad
 (a+b\varphi)^5\equiv a+3b\pmod{5R}.
\]
Thus a fifth power is a rational scalar modulo \(5R\).  The coordinate
pairs of \(1,\varphi,\ldots,\varphi^4\) are
\[
 (1,0),\quad(0,1),\quad(1,1),\quad(1,2),\quad(2,3).
\]
Suppose \(\alpha=A+B\varphi\), \(5\mid B\), \(5\nmid A\), and
\(\alpha\bar\alpha=c^5\), where \(\alpha,\bar\alpha\) are coprime in
\(R\).  Unique factorization shows that each prime multiplicity in
\(\alpha\) is a multiple of \(5\), since the prime occurs in no factor
of \(\bar\alpha\).  Hence
\(\alpha=\varepsilon\delta^5=\varphi^r\gamma^5\).
Modulo \(5R\), \(\gamma^5\) is a nonzero scalar: a zero scalar
would make both coordinates of \(\alpha\) divisible by \(5\).
The displayed table and \(5\mid B\) now force \(r=0\).
We have proved
\[
 \boxed{\quad
 \alpha\bar\alpha=c^5,\quad
 (\alpha,\bar\alpha)=1,\quad
 5\mid B,\quad5\nmid A
 \quad\Longrightarrow\quad \alpha=\gamma^5.
 \quad}
 \tag{Q1}\label{eq:quintic-extraction}
\]

Set
\[
 F(t,s)=t^4+10t^2s^2+5s^4,\qquad
 G(t,s)=t^4+50t^2s^2+125s^4.
\]
Expansion gives
\[
 (t+s\sqrt5)^5=tG(t,s)+5sF(t,s)\sqrt5.
 \tag{Q2}\label{eq:quintic-expansion}
\]
We record the two precise forms of (Q1) that will be used.

\begin{lemma}[Direct coordinates]\label{lem:quintic-direct}
Let \(P,Q,c\) be positive integers with
\[
 (P,Q)=1,\quad P\text{ odd},\quad Q\text{ even},\quad
 5\mid Q,\quad5\nmid P,\quad P^2-5Q^2=c^5.
\]
Then there are positive coprime integers \(t,s\), with \(t\) odd,
\(s\) even and \(5\nmid t\), such that
\[
 P=tG(t,s),\qquad Q=5sF(t,s).
 \tag{Q3}\label{eq:quintic-direct}
\]
\end{lemma}

\begin{proof}
Put \(\alpha=P+Q\sqrt5=(P-Q)+2Q\varphi\).
The integer \(\gcd(P^2-5Q^2,2P)\) is \(1\):
an odd prime common divisor would divide \(P\) and \(5Q^2\),
contrary to \((P,Q)=1\) and \(5\nmid P\), and the norm is odd.
A Bezout identity between the product
\(\alpha\bar\alpha=P^2-5Q^2\) and the sum
\(\alpha+\bar\alpha=2P\) gives a Bezout identity between
\(\alpha\) and \(\bar\alpha\).  Indeed
\(u\alpha\bar\alpha+v(\alpha+\bar\alpha)=1\) becomes
\((u\bar\alpha+v)\alpha+v\bar\alpha=1\).
The coordinate hypotheses of (Q1) hold, so \(\alpha=\gamma^5\).

Write \(\gamma=a+b\varphi\).  Modulo \(2R\) one has
\(\gamma^4=\gamma\), by checking the four coordinate pairs, and hence
\(\gamma^5=\gamma^2\).  The second coordinate of
\(\gamma^2\) is \(b^2\equiv b\pmod2\).
The second coordinate of \(\alpha\) is \(2Q\), so \(b\) is even.
Thus \(\gamma=t+s\sqrt5\) with integers
\(s=b/2\), \(t=a+b/2\).
Its norm is odd, so \(t,s\) have opposite parity.

Formula (Q2) proves (Q3).  Its two coefficient polynomials have positive
even-power terms.  Since \(P,Q>0\), both \(t,s\) are positive.
If \(t\) were even and \(s\) odd, \(5sF(t,s)\) would be odd,
contrary to the parity of \(Q\).  Thus \(t\) is odd and \(s\) even.
The first equation modulo \(5\) gives \(5\nmid t\).

Finally, \(\gamma,\bar\gamma\) are coprime: a Bezout identity between
their fifth powers is also a Bezout identity between them after
absorbing fourth powers into the coefficients.  If a rational prime
divided both \(t,s\), it would divide both \(\gamma,\bar\gamma\);
that prime is a nonunit of \(R\), since its norm is its square.
This is impossible.  Hence \((t,s)=1\).
\end{proof}

\begin{lemma}[Half-integral coordinates]\label{lem:quintic-half}
Let \(P,Q,c\) be positive integers with
\[
 (P,Q)=1,\quad P,Q\text{ odd},\quad
 5\mid Q,\quad5\nmid P,\quad P^2-5Q^2=4c^5.
\]
Then there are positive coprime odd integers \(t,s\), with \(5\nmid t\),
such that
\[
 16P=tG(t,s),\qquad16Q=5sF(t,s).
 \tag{Q4}\label{eq:quintic-half}
\]
\end{lemma}

\begin{proof}
Take
\(\alpha=(P+Q\sqrt5)/2=(P-Q)/2+Q\varphi\).
Its norm is \(c^5\), and its trace is \(P\).
The gcd of these two integers is \(1\), because a prime dividing both
would divide \(P\) and \(5Q^2\); also \(P\) is odd.
The product-and-sum Bezout argument again proves that
\(\alpha,\bar\alpha\) are coprime.  Its second coordinate is divisible
by \(5\) and its first is not.  Thus (Q1) gives \(\alpha=\gamma^5\).

Write \(\gamma=a+b\varphi=(t+s\sqrt5)/2\), where
\(s=b\), \(t=2a+b\).  The modulo-two calculation in the preceding
proof shows that \(b\) is odd, since the second coordinate of
\(\alpha\) is \(Q\).  Hence \(t,s\) are odd.
Multiplying \(\alpha=\gamma^5\) by \(32\) and applying (Q2) gives
(Q4).  Positivity of \(P,Q\) gives positivity of \(t,s\), and the
first coefficient modulo \(5\) gives \(5\nmid t\).
If an odd prime divided both \(t,s\), then \(\gamma\) and
\(\bar\gamma\) would both be divisible by that prime in \(R\);
the quotients remain in \(R\), because division of both odd
coordinates by an odd integer preserves their common parity.
Coprimality of the conjugates excludes this.  The prime \(2\) is
excluded because \(t,s\) are odd.
\end{proof}

The witnesses in these two lemmas can, if desired, be selected directly
by searching
\[
 1\le t,s\le16(P+Q+1).
 \tag{Q5}\label{eq:quintic-root-bound}
\]
The coefficient equations just proved place the witnesses inside this
box.  Thus no unbounded search for an algebraic fifth root occurs in
the actual descent.

\subsection{Two impossible families}

For \(h\ge0\), define an \emph{odd state} to be positive integers
\(t,s,w\) satisfying
\[
 \begin{gathered}
 (t,s)=1,\quad t,s\text{ odd},\quad5\mid s,\quad5\nmid t,\\
 5^h sF(t,s)=16w^5.
 \end{gathered}
 \tag{Q6}\label{eq:quintic-oddstate}
\]
Define an \emph{even state} by
\[
 \begin{gathered}
 (t,s)=1,\quad t\text{ odd},\quad s\text{ even},
 \quad5\mid s,\quad5\nmid t,\\
 2\cdot5^h sF(t,s)=w^5.
 \end{gathered}
 \tag{Q7}\label{eq:quintic-evenstate}
\]
The coefficient \(2\) in the even family stays fixed throughout the
descent; no additional parameter for its exponent is necessary.

The gcd calculations used below are
\[
 (s,F(t,s))=1,\qquad (5,F(t,s))=1.
 \tag{Q8}\label{eq:quintic-gcd}
\]
Indeed \(F(t,s)\equiv t^4\pmod s\) and
\(F(t,s)\equiv t^4\pmod5\).
For odd \(t,s\), \(16\mid F(t,s)\), since odd fourth powers are
\(1\pmod {16}\) and odd squares are \(1\) or \(9\pmod {16}\).
If \(t\) is odd and \(s\) even, \(F(t,s)\) is odd.

In an odd state, divide \(F\) by \(16\) and use (Q8).
The two coprime positive factors
\(5^h s\) and \(F(t,s)/16\) have product \(w^5\); therefore
\[
 5^h s=u^5,\qquad F(t,s)=16v^5
 \tag{Q9}\label{eq:quintic-oddsplit}
\]
for positive integers \(u,v\).
In an even state the corresponding conclusion is
\[
 2\cdot5^h s=u^5,\qquad F(t,s)=v^5.
 \tag{Q10}\label{eq:quintic-evensplit}
\]
These are ordinary integer prime-multiplicity arguments, with bounded
factorization or bounded searches for \(u,v\).

Now use the identity
\[
 (t^2+5s^2)^2-5(2s^2)^2=F(t,s).
 \tag{Q11}\label{eq:quintic-normidentity}
\]
For an odd state set \(P=(t^2+5s^2)/2\), \(Q=s^2\).
Both are odd, their gcd is \(1\), \(5\mid Q\), \(5\nmid P\), and
\(P^2-5Q^2=4v^5\).
Here \((P,s^2)=1\) follows from \((t,s)=1\), and the other claims
follow directly by reduction modulo \(4\) and \(5\).
Lemma~\ref{lem:quintic-half} produces positive coprime odd \(t',s'\)
with \(5\nmid t'\) and
\[
 16s^2=5s'F(t',s').
 \tag{Q12}\label{eq:quintic-oddnext}
\]
For an even state set \(P=t^2+5s^2\), \(Q=2s^2\).
Then \(P\) is odd, \(Q\) even, they are coprime, \(5\mid Q\),
\(5\nmid P\), and \(P^2-5Q^2=v^5\).
Lemma~\ref{lem:quintic-direct} gives positive coprime \(t',s'\),
\(t'\) odd, \(s'\) even, \(5\nmid t'\), and
\[
 2s^2=5s'F(t',s').
 \tag{Q13}\label{eq:quintic-evennext}
\]
In either case \(5\mid s'\).  The left side is divisible by \(25\);
cancelling one factor \(5\) on the right shows
\(5\mid s'F(t',s')\), while \(5\nmid F(t',s')\).

The new second coordinate is strictly smaller:
\[
 \boxed{s'<s.}
 \tag{Q14}\label{eq:quintic-decrease}
\]
In fact \(F(t',s')>5s'^4\), because \(t'>0\).
Writing either (Q12) or (Q13) as
\(Ks^2=5s'F(t',s')\), with \(K=16\) or \(2\), gives
\(25s'^5<Ks^2\).
If \(s'\ge s\), positivity gives
\(s^2\le s'^2\le s'^5\), contradicting \(K<25\).

Finally, square the first equation in (Q9) or (Q10).
For the odd state, (Q12) gives
\[
 5^{2h+1}s'F(t',s')
 =16(5^h s)^2
 =16(u^2)^5.
\]
For the even state, (Q13) gives
\[
 2\cdot5^{2h+1}s'F(t',s')
 =4\cdot5^{2h}s^2
 =(u^2)^5.
\]
Thus the successor is a state of the same type with
\[
 (h,t,s,w)\longmapsto(2h+1,t',s',u^2)
 \quad\hbox{and}\quad0<s'<s.
 \tag{Q15}\label{eq:quintic-transition}
\]

Neither state family has a member.  To see this in precisely the form
needed for iPRA, totalize the successor by returning a fixed default
on invalid input, and otherwise use the bounded searches described
above.  Starting from a purported state with second coordinate \(s_0\),
iterate this primitive-recursive function \(s_0\) times.
Quantifier-free induction on the iteration number proves that each
iterate remains a state and that
\(s_i+i\le s_0\).
At \(i=s_0\) this gives \(s_i\le0\), whereas every state has
\(s_i>0\).  This is a finite contradiction, not an invocation of
a well-ordering principle for arbitrary predicates.

\subsection{A fifth-power Fermat equation enters one of the families}

Suppose nonzero integers satisfy \(a^5+b^5=c^5\).
Cancelling their common gcd gives a pairwise coprime solution.

\paragraph{The modulo-25 preliminary.}
A fifth power modulo \(25\) is one of
\(0,1,7,-7,-1\); the nonzero four residues correspond exactly to
bases not divisible by \(5\).
No sum of three members of \(\{1,7,-7,-1\}\) is \(0\pmod {25}\).
For instance, their pairwise sums are
\(0,\pm2,\pm6,\pm8,\pm14\), none of which is
\(1,7,-7,-1\pmod {25}\).
Applying this to \(a^5+b^5+(-c)^5=0\) shows that one base is divisible
by \(5\).

Move that term to the right, changing signs if necessary, to obtain
\[
 x^5+y^5=5^5z^5,\qquad
 (x,y)=1,\qquad xyz\ne0,\qquad z>0.
 \tag{Q16}\label{eq:quintic-normalized}
\]
Neither \(x\) nor \(y\) is divisible by \(5\).
They cannot both be even.

\paragraph{Opposite parity.}
If \(x,y\) have opposite parity, put
\(p=x+y\), \(q=x-y\).  Then \(p,q\) are odd,
\((p,q)=1\), and
\[
 p(p^4+10p^2q^2+5q^4)=16\cdot5^5z^5.
\]
The quartic factor is positive, so \(p>0\).
Modulo \(5\) the left side is \(p^5\), hence \(p=5r\) with \(r>0\).
Coprimality implies \(5\nmid q\) and \((q,r)=1\);
in particular \(q\ne0\).
Writing \(q_0=|q|\), we get positive coprime odd \(q_0,r\), with
\(5\nmid q_0\), and
\[
 rG(q_0,r)=16\cdot5^3z^5.
 \tag{Q17}\label{eq:quintic-core}
\]

\paragraph{Both odd.}
If \(x,y\) are odd, put
\(p=(x+y)/2\), \(q=(x-y)/2\).
These are coprime and have opposite parity, and
\[
 2p(p^4+10p^2q^2+5q^4)=5^5z^5.
\]
Again \(p>0\) and \(p=5r\) with \(r>0\).
The right side is even, so \(z=2z_0\) with \(z_0>0\).
Substitution and cancellation give exactly
\[
 rG(|q|,r)=16\cdot5^3z_0^5.
\]
If \(r\) were odd, \(q\) would be even, and the left side would be
odd, a contradiction.  Thus \(r\) is even and \(|q|\) odd.
Again \((|q|,r)=1\), \(5\nmid q\), and \(q\ne0\).

We have therefore reached the common core
\[
 rG(q,r)=16\cdot5^3Z^5,\qquad
 q,r,Z>0,\quad (q,r)=1,\quad q\text{ odd},\quad5\nmid q,
 \tag{Q18}\label{eq:quintic-unified-core}
\]
where \(r\) is either odd or even.
We now spell out the last allocation of coprime factors.

The factors \(r,G(q,r)\) are coprime, since
\(G(q,r)\equiv q^4\pmod r\), and \(5\nmid G(q,r)\).
If \(r\) is odd, \(16\mid G(q,r)\); this follows from
odd fourth powers modulo \(16\) and
\(50\equiv2,\ 125\equiv13\pmod {16}\).
Thus
\[
 (25r)\frac{G(q,r)}{16}=(5Z)^5
\]
is a product of coprime positive integers, and so
\[
 25r=A^5,\qquad G(q,r)=16B^5.
 \tag{Q19}\label{eq:quintic-initial-odd-split}
\]
If \(r\) is even, \(G(q,r)\) is odd; hence
\[
 (50r)G(q,r)=(10Z)^5
\]
is a coprime factorization and
\[
 50r=A^5,\qquad G(q,r)=B^5.
 \tag{Q20}\label{eq:quintic-initial-even-split}
\]
Both cases imply \(5\mid r\): in the first,
\(5\mid A\) and \(r=125(A/5)^5\); in the second,
\(2r=125(A/5)^5\), which has the same consequence.

Use the second norm identity
\[
 G(q,r)=(q^2+25r^2)^2-5(10r^2)^2.
 \tag{Q21}\label{eq:quintic-initial-norm}
\]
For odd \(r\), take
\[
 P=(q^2+25r^2)/2,\qquad Q=5r^2.
\]
Then \(P,Q\) are odd, positive, and coprime,
\(5\mid Q\), \(5\nmid P\), and
\(P^2-5Q^2=4B^5\).
The coprimality is immediate from \((q,r)=1\), \(5\nmid q\),
and the oddness of \(P\).
Lemma~\ref{lem:quintic-half} gives positive coprime odd \(t,s\),
\(5\nmid t\), with
\[
 16r^2=sF(t,s).
\]
Since \(5\mid r\) and \(F(t,s)\equiv t^4\pmod5\), we have \(5\mid s\).
Squaring \(25r=A^5\) proves
\[
 5^4sF(t,s)=16(A^2)^5.
\]
This is an odd state with \(h=4\).

For even \(r\), take
\(P=q^2+25r^2\), \(Q=10r^2\).
These are coprime with \(P\) odd and \(Q\) even;
\(5\mid Q\), \(5\nmid P\), and
\(P^2-5Q^2=B^5\).
Lemma~\ref{lem:quintic-direct} gives positive coprime \(t,s\),
\(t\) odd, \(s\) even, \(5\nmid t\), and
\[
 2r^2=sF(t,s).
\]
Again \(5\mid s\).  Squaring \(50r=A^5\) yields
\[
 2\cdot5^4sF(t,s)=(A^2)^5,
\]
an even state with \(h=4\).

Both families were proved impossible.  This contradiction establishes:
\[
 \boxed{\forall a,b,c>0,\qquad a^5+b^5\ne c^5.}
 \tag{Q22}\label{eq:flt-five}
\]

\subsection{Primitive-recursive and Lean audit for the quintic proof}

All searches in the proof have explicit finite bounds.
Integer factorization can use trial division through the input.
The Euclidean algorithm in \(R\) has at most one more iteration than
the initial absolute norm.  Normalization of an element with coordinate
height \(H\) searches signed exponents between \(-4H\) and \(4H\).
Factorization in \(R\) searches divisor coordinates bounded by \(3M+6\),
where \(M\) is its absolute norm.  Algebraic-root extraction used by
the descent can instead search the box (Q5).
A fifth root of a positive integer \(D\) is sought only in \([1,D]\).
The main descent is iterated exactly \(s_0\) times.

For an especially simple envelope on the main states, if their four
coordinates \(h,t,s,w\) are at most \(B\ge1\), the next coordinates
are bounded by
\[
 T(B)=2^{100(B+1)}.
\]
Indeed \(h'=2h+1\), the two norm coordinates are bounded by
\(6(B+1)^2\), (Q5) bounds \(t',s'\) polynomially in \(B\), and
\(w'=u^2\le(2\cdot5^h s)^2\).
Iterating \(T\) exactly \(s_0\) times supplies one primitive-recursive
bound for the complete main-state trace.
For the internal Euclidean calculations, one may bound a coordinate
height after each division by
\(E(H)=32(H+1)^3\), and iterate \(E\) for the already bounded
number of divisions.  Efficiency is irrelevant; these bounds certify
primitive recursiveness.

A Lean completion should use integer-pair code for \(R\), not
\texttt{Real.sqrt}, and should isolate the following interfaces:
\[
\begin{array}{ll}
\texttt{goldenEuclideanStep}&
  \text{explicit quotient, remainder, and norm decrease},\\
\texttt{goldenUnitNormalForm}&
  u=\pm\varphi^j\text{ with a bounded signed exponent},\\
\texttt{goldenFifthExtraction}&
  \text{the coprime extraction implication (Q1)},\\
\texttt{fifthCoordinatesDirect}&
  \text{the direct bounded coordinate lemma},\\
\texttt{fifthCoordinatesHalf}&
  \text{the half-integral coordinate lemma},\\
\texttt{fifthOddStep},\ \texttt{fifthEvenStep}&
  \text{the transitions (Q15)},\\
\texttt{fifthInitialState}&
  \text{the finite normalization (Q16)--(Q21)}.
\end{array}
\]
The state predicates are decidable combinations of integer equalities,
inequalities, gcd, parity, and remainders.
Each interface should be proved as a quantifier-free correctness
statement about a specified primitive-recursive selector.
A final quantifier-free invariant on the bounded trace gives the iPRA
derivation.  Proving the same mathematical implication in Lean's
ambient logic is a separate, weaker deliverable than exhibiting that
object-theoretic iPRA derivation.

For orientation, the independently available source tree
\texttt{isomorphic-ai/fermat}, under
\texttt{Fermat/Five/}, separates the historical argument into
\texttt{Coordinates}, \texttt{PowerExtraction}, \texttt{Initial}, and
\texttt{Descent} modules.  It was inspected as a comparison source,
not compiled as part of this package.  Its ambient use of real embeddings,
choice, or least-counterexample formulations must not be counted as
an iPRA formalization.  The explicit bounded normalization and fixed-length
iteration above describe the replacements required here.


\section{The seventh exponent: Lebesgue's descent, with the missing case repaired}
\label{sec:seven}

This section gives a complete elementary proof for exponent $7$.  The algebraic
route is due to Lebesgue, including his published correction
\cite{Lebesgue1840,Lebesgue1840Add}.  We spell out the factor allocations that
make the correction indispensable.  All factorizations below are factorizations
of ordinary integers; no number field, ideal class group, or algebraic number
theory is used.

We use the following elementary finite-arithmetic conventions.  For a prime
$\ell$ and a nonzero integer $m$, $v_\ell(m)$ is the largest $j\leq |m|$ for
which $\ell^j\mid m$.  The value is obtained by bounded search.  Trial division
gives prime factorizations of positive integers, and if two positive coprime
factors have product $w^k$, each is a $k$th power: each prime occurs in exactly
one factor and its exponent is a multiple of $k$.  These assertions admit
bounded proofs by scanning integers up to the integer being factored.  Signs
are treated separately.

\subsection{An elementary quartic descent}

\begin{lemma}[Lebesgue's quartic obstruction]
\label{lem:seven-quartic}
There are no integers $a,P,Q,R$ such that $a\geq1$, $P,Q,R$ are positive odd
and pairwise coprime, and
\begin{equation}\label{eq:seven-quartic}
 P^2=Q^4-3\cdot7^4\,2^{2a}Q^2R^2
              +7^7\,2^{4a+4}R^4.
\end{equation}
The same conclusion holds when $P,Q,R$ are nonzero signed odd integers, by
taking absolute values.
\end{lemma}

\begin{proof}
Write $D_a(P,Q,R)$ for the displayed conditions.  If $a=1$, the right side of
\eqref{eq:seven-quartic} is $1-4\equiv5\pmod8$, whereas the square of an odd
integer is $1\pmod8$.  Hence $a\geq2$.

Put
\[
 C=Q^2-3\cdot7^4\,2^{2a-1}R^2.
\]
Completing the square gives
\begin{equation}\label{eq:seven-first-product}
 (P+C)(P-C)=7^7\,2^{4a-2}R^4.
\end{equation}
Both factors are positive, since $P^2-C^2>0$ and $P>0$.  They have greatest
common divisor $2$.  To check this last assertion, first $7\nmid Q$: otherwise
the original equation gives $7\mid P$, contrary to $\gcd(P,Q)=1$.  Any odd
prime dividing both factors would divide $7R$, by their product, and would
also divide $C$.  A prime dividing $R$ cannot divide $C$, since
$C\equiv Q^2$ modulo that prime and $\gcd(Q,R)=1$.  The prime $7$ cannot divide
$C$ either.  Finally, $P,C$ are odd and $(P+C)+(P-C)=2P$ has exactly one factor
of $2$.  Thus one factor in \eqref{eq:seven-first-product} has $2$-valuation
$1$, the other $4a-3$.

Unique factorization in the integers now gives positive odd coprime $s,t$,
with $R=st$, in one of the following four allocations:
\begin{center}
\begin{tabular}{c|cc}
 & $P+C$ & $P-C$ \\
\hline
I   & $2s^4$ & $2^{4a-3}7^7t^4$ \\
II  & $2\cdot7^7s^4$ & $2^{4a-3}t^4$ \\
III & $2^{4a-3}s^4$ & $2\cdot7^7t^4$ \\
IV  & $2^{4a-3}7^7s^4$ & $2t^4$
\end{tabular}
\end{center}
These four rows include the possibility that $7$ already divides $R$:
the factor receiving $7^7$ then receives the entire additional fourth power
of $7$ from $R^4$.  The factors in the other column are coprime to $7$.

Since $C\equiv1\pmod4$, rows II and IV are impossible: subtracting the
columns and dividing by $2$ gives $C\equiv3\pmod4$ in both.  Row III is
\emph{not} ruled out by this congruence.  We first treat row I, then eliminate
row III separately.

\paragraph{Row I produces a smaller descent parameter.}
Subtracting its two equations gives
\[
 Q^2=s^4+3\cdot7^4\,2^{2a-1}s^2t^2
                    -7^7\,2^{4a-4}t^4.
\]
Let $B=s^2+3\cdot7^4\,2^{2a-2}t^2$.  Direct multiplication gives
\begin{equation}\label{eq:seven-second-product}
 (B+Q)(B-Q)=7^7\,2^{4a+2}t^4.
\end{equation}
Here $B>|Q|$ and both factors are positive.  Their gcd is $2$: an odd common
prime would divide $7t$ and $B$, but modulo such a prime $B$ is a nonzero
multiple of $s^2$.  For $7$ we use $7\nmid s$, which follows from row I
and the gcd-$2$ assertion for its two factors.  For a divisor of $t$ we use
$\gcd(s,t)=1$.  Exactly one of $B+Q,B-Q$ has $2$-valuation $1$; we may replace
$Q$ by $-Q$ to make it the first.  Consequently, for positive odd coprime
$u,v$ with $t=uv$, either
\[
 B+Q=2u^4,\qquad B-Q=2^{4a+1}7^7v^4,
\]
or
\[
 B+Q=2\cdot7^7u^4,\qquad B-Q=2^{4a+1}v^4.
\]
The second possibility gives $B\equiv3\pmod4$, whereas the definition of
$B$ gives $B\equiv1\pmod4$.  The first possibility therefore holds.  Adding
its two equations yields
\[
 s^2=u^4-3\cdot7^4\,2^{2a-2}u^2v^2
                    +7^7\,2^{4a}v^4.
\]
This is exactly $D_{a-1}(s,u,v)$.  All three new integers are positive and
odd.  They are pairwise coprime because $R=st=suv$, $\gcd(s,t)=1$, and
$\gcd(u,v)=1$.  Moreover
\begin{equation}\label{eq:seven-descent-bound}
 s,u,v\leq R.
\end{equation}

\paragraph{Row III is impossible: the correction.}
In this row subtraction instead gives
\[
 Q^2=2^{4a-4}s^4+3\cdot7^4\,2^{2a-1}s^2t^2-7^7t^4.
\]
Set $E=2^{2a-2}s^2+3\cdot7^4t^2$.  Then
\begin{equation}\label{eq:seven-correction-product}
 (E+Q)(E-Q)=2^6\,7^7t^4.
\end{equation}
Again both factors are positive and their gcd is $2$.  The proof is the
same explicit prime-divisor argument: an odd common prime divides $7t$;
$7\nmid s$ follows from row III; and $\gcd(s,t)=1$.
Choose the sign of $Q$ so that $E+Q$ has $2$-valuation $1$.  With $t=uv$ and
$u,v$ positive odd coprime, the two possible allocations are
\begin{align*}
 E+Q=2v^4,&\qquad E-Q=2^5\,7^7u^4,\\
 \text{or}\qquad E+Q=2\cdot7^7v^4,&\qquad E-Q=2^5u^4.
\end{align*}
The first gives
\[
 2^{2a-2}s^2=v^4-3\cdot7^4u^2v^2+2^4\,7^7u^4,
\]
whose left side is $0\pmod4$ and right side $2\pmod4$.  Hence only the
second allocation remains, giving
\[
 2^{2a-2}s^2=7^7v^4-3\cdot7^4u^2v^2+2^4u^4.
\]
Multiplying by $2^6$ and completing the square once more gives
\begin{equation}\label{eq:seven-final-correction}
 (2^{a+2}s)^2-(2^5u^2-3\cdot7^4v^2)^2=7^7v^4.
\end{equation}
Put $F=2^{a+2}s$ and $G=2^5u^2-3\cdot7^4v^2$.  The factors $F+G,F-G$
are positive odd and coprime.  Indeed, any common odd prime divides $7v$;
if it divides $v$, the expression for $G$ and $\gcd(u,v)=1$ exclude it;
if it is $7$, the second allocation above shows $7\nmid u$, also excluding
it.  Integer factorization therefore gives positive odd $m,n$ with $v=mn$
such that, in some order, $F+G=m^4$ and $F-G=7^7n^4$.  Adding yields
\[
 2^{a+3}s=m^4+7^7n^4.
\]
Every odd fourth power is $1\pmod{16}$ and $7^7\equiv7\pmod{16}$.
The right side is therefore $8\pmod{16}$, while the left side is
$0\pmod{16}$.  This contradiction eliminates row III.

Thus every solution with $a\geq2$ produces a solution with parameter $a-1$,
and parameter $1$ is impossible.  The iteration has at most $a-1$ steps.
It is a bounded descent, and proves the lemma.
\end{proof}

\subsection{Reducing a seventh-power equation to the quartic}

\begin{theorem}\label{thm:flt-seven}
For positive integers $A,B,C$ one has $A^7+B^7\ne C^7$.
\end{theorem}

\begin{proof}
Suppose $A^7+B^7=C^7$.  Divide by the gcd to make the triple primitive;
prime divisibility shows that it is then pairwise coprime.  Put
$x=A$, $y=B$, $z=-C$, so $x^7+y^7+z^7=0$, and define
\begin{align*}
 s&=x+y+z, &
 u&=x^2+y^2+z^2+xy+xz+yz,\\
 v&=(x+y)(y+z)(z+x), &
 t&=u^2+xyzs.
\end{align*}
The elementary symmetric identities we require are
\begin{equation}\label{eq:seven-symmetric-identities}
 s^7=7vt,\qquad
 v=s(xy+xz+yz)-xyz,\qquad
 u=s^2-(xy+xz+yz).
\end{equation}
The first follows by expanding
$(x+y+z)^7-(x^7+y^7+z^7)$; its factorization is
$7(x+y)(y+z)(z+x)(u^2+xyzs)$.  The other two follow by multiplication.

There is exactly one even member of $x,y,z$: there cannot be two in a
primitive triple, and the sum of three odd seventh powers is odd.  Hence
$s,v$ are even, $u$ is odd, and $t\equiv1\pmod4$.
Furthermore $C>A,B$ and $C<A+B$.  The latter follows from
$(A+B)^7>A^7+B^7$ and the strict monotonicity of the seventh power.
Consequently $s>0$, $v=(A+B)(B-C)(A-C)>0$, and the first identity gives
$t>0$.  Also $u>0$, since
$2u=s^2+x^2+y^2+z^2$.

\paragraph{The coprimality information.}
We claim
\begin{equation}\label{eq:seven-coprimality}
 \gcd(t,xyz)=1,\qquad \gcd(t,v)=1.
\end{equation}
If a prime $\ell$ divides $t$, then $\ell\mid s$ by $s^7=7vt$, and
$\ell\mid u$ by $t=u^2+xyzs$.  If it also divides, say, $x$, then
$z\equiv-y\pmod\ell$, and $u\equiv y^2\pmod\ell$ gives $\ell\mid y$,
contrary to pairwise coprimality.  This proves the first assertion.  If
$\ell$ divides both $t$ and $v$, then $\ell\mid s$, and the middle identity
in \eqref{eq:seven-symmetric-identities} gives $\ell\mid xyz$, contradicting
the first assertion.  This proves the second.

\paragraph{Why $t$ is a fourteenth power.}
Fix a prime $\ell\mid t$.  Write
\[
 d=v_\ell(s),\qquad e=v_\ell(u),\qquad f=v_\ell(t),\qquad
 \epsilon=\begin{cases}1&\ell=7,\\0&\ell\ne7.\end{cases}
\]
Here $d,e,f>0$.  As $\ell\nmid v$, the first identity gives
$f=7d-\epsilon>d$.  As $\ell\nmid xyz$, the terms in
$t=u^2+xyzs$ have valuations $2e$ and $d$.  If these valuations were
unequal, the valuation of their sum would be their minimum, which is at
most $d$, a contradiction.  Thus $d=2e$, and
\begin{equation}\label{eq:seven-valuations}
 f=14e-\epsilon.
\end{equation}
For completeness, the unequal-valuation rule is obtained by extracting
the smaller power of $\ell$: the remaining summand of lower valuation is
nonzero modulo $\ell$, whereas the other summand is zero modulo $\ell$.

If $7\nmid t$, equation \eqref{eq:seven-valuations} says that every prime
exponent in $t$ is a multiple of $14$.  If $7\mid t$, it instead says that
every prime exponent in $7t$ is a multiple of $14$.  The latter alternative
is impossible, since $7t\equiv3\pmod4$ is not a square.  We have therefore
proved
\[
 t=q^{14},\qquad 7\nmid q
\]
for a positive odd integer $q$.  Moreover, for every prime dividing $q$,
equation \eqref{eq:seven-valuations} shows that its exponent in $u$ is
exactly its exponent in $q$.  We may write
\[
 u=qr,\qquad \gcd(q,r)=1,
\]
with $r$ positive odd.  The assertion includes $q=1$, where it is immediate.
The same valuation calculation gives $q^2\mid s$.  On writing $s=q^2k$,
the identity $s^7=7vq^{14}$ becomes $k^7=7v$.  Thus $7\mid k$, and for a
positive integer $h$ we obtain
\begin{equation}\label{eq:seven-parameters}
 s=7hq^2,\qquad v=7^6h^7,\qquad u=qr,\qquad t=q^{14}.
\end{equation}
The integer $h$ is even.  Also $\gcd(h,q)=1$, since $\gcd(t,v)=1$.

We need one more gcd: $\gcd(h,r)=1$.  Suppose a prime $\ell$ divided both.
Then it would divide $s,u,v$, by \eqref{eq:seven-parameters}, and would
divide $xyz$ by $v=s(xy+xz+yz)-xyz$.  If it divided $x$, then
$z\equiv-y\pmod\ell$ and $u\equiv y^2\pmod\ell$, forcing $\ell\mid y$.
The other choices are identical.  This contradicts primitivity.

\paragraph{The final square completion.}
Using $xy+xz+yz=s^2-u$ and $xyz=s(s^2-u)-v$ in $t=u^2+xyzs$, and then
substituting \eqref{eq:seven-parameters} and dividing by $q^2$, gives
\[
 q^{12}=r^2+7^2h^2q^3(7^2h^2q^3-r)-7^7h^8.
\]
Consequently
\begin{equation}\label{eq:seven-reduction-square}
 \left(r-\frac{7^2h^2q^3}{2}\right)^2
 =q^{12}-3\left(\frac{7^2h^2q^3}{2}\right)^2+7^7h^8.
\end{equation}
Extract the power of $2$ from $h^2$, writing $h^2=2^{a+1}R$ with $R$ odd.
Since $h$ is even, $a\geq1$.  Define
\[
 P=\left|r-\frac{7^2h^2q^3}{2}\right|,\qquad Q=q^3.
\]
The integer $P$ is positive odd: $r$ is odd and $h^2/2$ is even.
The three integers $P,Q,R$ are pairwise coprime.  Indeed $\gcd(Q,R)=1$
follows from $\gcd(q,h)=1$; reducing $P$ modulo any prime divisor of $Q$
reduces it to $\pm r$; and reducing it modulo any prime divisor of $R$
does the same.  We use $\gcd(q,r)=\gcd(h,r)=1$ in these last two cases.
Finally, \eqref{eq:seven-reduction-square} is exactly
\[
 P^2=Q^4-3\cdot7^4\,2^{2a}Q^2R^2+7^7\,2^{4a+4}R^4.
\]
This contradicts Lemma~\ref{lem:seven-quartic}.
\end{proof}

\subsection{The iPRA audit and Lean implementation path for exponent 7}

Every search used above has a primitive-recursive bound.  Integer powers,
signed arithmetic, gcd, prime testing, valuations of nonzero integers,
factorization by trial division, and extraction of a verified exact root
are primitive recursive.  An exact $k$th root of $N>0$ is searched for only
in $[1,N]$; its existence is proved by the prime-exponent calculation, not
assumed from the search.  The prime-exponent calculation itself scans a
finite list of divisors of $N$.

There is also a simple bound in terms of the original input alone.  Put
$N=\max(A,B,C)$.  Primitive reduction does not increase $N$; the displayed
definitions give $s\leq3N$ and $u\leq6N^2$.  From
$s=7hq^2$ and positivity, both $h,q\leq3N$, while $r\leq6N^2$.
Consequently the final quartic has
\[
 a\leq h^2\leq9N^2,\qquad
 P,Q,R\leq M:=2^{14}(N+1)^5.
\]
For example $Q=q^3\leq27N^3$, $R\leq h^2\leq9N^2$, and
$P\leq r+49h^2q^3/2<2^{14}(N+1)^5$.  The outer descent therefore has
at most $9N^2$ steps, and all its coordinates remain bounded by $M$.
The larger integers temporarily used in the square completions have
explicit bounds obtained by substituting these bounds into their displayed
polynomials and powers of $2$.  Such bounds are primitive recursive.

The logical point in the quartic descent merits explicit treatment.  Fix a
bound $M$ and let
\[
 \operatorname{Bad}_7(a,M)=1
\]
mean that the bounded scan of $1\leq P,Q,R\leq M$ finds a triple satisfying
$D_a(P,Q,R)$.  This is a primitive-recursive Boolean function.  The proof
above gives $\operatorname{Bad}_7(1,M)=0$ and, by
\eqref{eq:seven-descent-bound},
\[
 \operatorname{Bad}_7(a,M)=0\ \Longrightarrow\
 \operatorname{Bad}_7(a+1,M)=0 \qquad(a\geq1).
\]
The induction formula is the equation
$\operatorname{Bad}_7(a+1,M)=0$, hence quantifier free.  The parameter $M$
is arbitrary but fixed during induction.  Given an alleged quartic
solution, take $M=\max(P,Q,R)$ and instantiate the result.  This is ordinary
iPRA induction; it is not induction on an unbounded existential assertion
and does not use a least counterexample principle.

All intermediate implications in this proof are intuitionistically valid.
Prime divisibility, comparisons of integer valuations, parity, and the four
factor allocations are decidable finite alternatives.  A contradiction is
derived from a supplied counterexample; no double-negation elimination for
an arbitrary formula is invoked.

For a transparent Lean development the suggested dependency order is:
\begin{enumerate}
\item Define signed-integer coding, the Boolean prime and coprime tests,
      bounded valuation, and bounded factorization.  Prove multiplication
      and unequal-sum valuation laws and the coprime power-splitting lemma.
\item Define the integer polynomial \texttt{quartic7 a P Q R}.  Check the
      four square-completion identities by ring normalization.  Store the
      four allocations as a finite inductive type, or four Boolean cases.
\item Prove \texttt{quartic7\_base}, \texttt{quartic7\_caseIII\_false}, and
      \texttt{quartic7\_step}.  The step returns an explicit triple
      $(s,u,v)$ and the three inequalities $s,u,v\leq R$.
\item Define the finite Boolean scan \texttt{bad7 a M} and prove its
      quantifier-free induction invariant.  This supplies the quartic
      obstruction internally to the chosen iPRA syntax.
\item Verify the symmetric polynomial identity and prove the valuation
      extraction lemma yielding \eqref{eq:seven-parameters}, including the
      two possible $7$-valuations and their modulo-$4$ separation.
\item Build the explicit transformation from a positive primitive
      seventh-power solution to $(a,P,Q,R)$, and compose it with the
      quartic obstruction.  Add the ordinary gcd reduction first or last.
\end{enumerate}
The equations proved by Lean's ambient kernel are a convenient intermediate
check.  A claim of an iPRA derivation additionally requires the finite scans,
induction formulas, and their proofs to be represented in the object theory;
the preceding construction specifies these data.

\section{The resulting infinite family and further specializations}
\label{sec:coverage}

\begin{lemma}[Exponent lifting]\label{lem:lifting}
For $d\geq1$ and $k\geq1$, $\ipra$ proves
\[
 \FLT(d)\longrightarrow\FLT(dk).
\]
\end{lemma}

\begin{proof}
Suppose $x,y,z>0$ and $x^{dk}+y^{dk}=z^{dk}$. The exponent law
$(x^k)^d=x^{dk}$ is proved by quantifier-free induction from the defining
recursion for powers and multiplication. The three positive integers
$x^k,y^k,z^k$ then contradict $\FLT(d)$. Positivity of these powers is
also proved by quantifier-free induction. All variables remain arbitrary,
so the implication is uniform in $d,k$.
\end{proof}

\begin{theorem}[Main exponent family]\label{thm:main}
The theory $\ipra$ proves the single uniform statement
\[
 \boxed{\quad
 \forall n\geq3\quad
 (3\mid n\ \vee\ 4\mid n\ \vee\ 5\mid n\ \vee\ 7\mid n)
 \ \longrightarrow\ \FLT(n).
 \quad}
\]
It also proves the equivalent nontrivial-solution exclusion over signed
integers and over coded rational numbers.
\end{theorem}

\begin{proof}
The four preceding proofs establish the four base exponents. In each
branch of the displayed disjunction, division produces $k=n/d$ with
$n=dk$ and $k\geq1$; apply Lemma~\ref{lem:lifting}.
For even $n$, taking absolute values reduces a signed equation to the
positive case. For odd $n$, write the signed equation as a sum of three
signed $n$th powers equal to zero. The three nonzero terms cannot have
the same sign. Move the term with the exceptional sign to the other side
to obtain a positive equation. These are finite decidable sign cases.
For nonzero rationals, choose the product of their three positive
denominators as a common denominator. Multiplication by its $n$th power
gives a nonzero signed integral solution, which has just been excluded.
All denominator clearing and sign tests are primitive recursive.
\end{proof}

This is an infinite collection of full Fermat specializations, including
all powers of $2$ of exponent at least two, all powers of $3,5,7$, and
every positive multiple of any of the four base exponents. For example,
up to $40$ it proves precisely
\[
 \begin{gathered}
 3,4,5,6,7,8,9,10,12,14,15,16,18,20,21,24,25,\
 27,28,30,32,33,35,36,39,40.
 \end{gathered}
\]
The first exponents $\geq3$ outside this package are
$11,13,17,19,22,23,26,29,31,34,37,38$.
This list describes the reach of the proof, not the truth value of FLT
at those exponents.

\paragraph{Exact coverage.}
The divisibility condition is periodic modulo $420=4\cdot3\cdot5\cdot7$.
The complement has
$3\cdot2\cdot4\cdot6=144$ residue classes: three allowed residues modulo
$4$, two modulo $3$, four modulo $5$, and six modulo $7$. The elementary
Chinese remainder bijection, or a direct scan of $0,\ldots,419$, therefore
gives $420-144=276$ covered classes. Its natural density is
\[
 \frac{276}{420}=\frac{23}{35}\approx65.7142857\%.
\]
There is a completely finite version of this calculation: in every full
block of $420$ consecutive integers exactly $276$ satisfy the divisibility
test. If $N=420q+r$, $0\leq r<420$, the count among $1,\ldots,N$
differs from $(23/35)N$ by less than $420$. Restricting to exponents
$n\geq3$ affects only the initial terms. No analytic limit principle is
needed for this counting estimate.

\begin{corollary}[Mixed exponents]\label{cor:mixed}
Let $a,b,c$ be positive integers. If one of $3,4,5,7$ divides all three
of $a,b,c$, then $\ipra$ proves that
\[
 X^a+Y^b=Z^c
\]
has no positive integral solution. Independently, for all $i,j,k\geq1$
it proves that
\[
 X^{4i}+Y^{4j}=Z^{2k}
\]
has no positive integral solution.
\end{corollary}

\begin{proof}
For the first assertion write $a=d\alpha$, $b=d\beta$, and
$c=d\gamma$ and apply $\FLT(d)$ to
$X^\alpha,Y^\beta,Z^\gamma$. For the second apply
Theorem~\ref{thm:four} to $X^i,Y^j,Z^k$.
\end{proof}

\paragraph{What has not been claimed.}
This paper does not supply a proof of general FLT in $\ipra$. It does
not establish exponent $11$, all regular prime exponents, or an
equivalence between ordinary Lean provability and provability in $\ipra$.
In particular, computing Bernoulli numerators or checking a regularity
condition would not by itself supply an $\ipra$ proof of Kummer's theorem.
We make no claim of independence or unprovability for the omitted cases.
The additional work needed for a new prime exponent is a complete
arithmetic argument with explicit constructions and induction invariants,
comparable to those given here.


\section{A concrete path to complete Lean formalization}
\label{sec:formalization}

\subsection{Three distinct levels of completion}

The mathematical proofs in this paper are complete human-readable
arguments, with their computation and induction bounds given explicitly.
The supplied Lean files are an implementation blueprint. They have not
been compiled in this environment, and this package contains no finished
Lean term proving the four base cases or their internal $\ipra$
derivations. This distinction is part of the deliverable, not an implicit
assumption about a future formalization.

There are three levels to keep separate:
\begin{enumerate}
\item \textbf{Arithmetic semantics.} Prove the displayed integer identities,
number-theoretic lemmas, and FLT assertions in Lean's ordinary logic.
\item \textbf{Programs and restricted induction.} Represent every numerical
operation by an explicit primitive recursive code, and identify every
induction formula as quantifier free. A termination proof alone is not
enough to establish this level.
\item \textbf{An internal derivation.} Construct an inhabitant of
\texttt{Prf PRAAx [] sentence} in ProveIt's actual syntax, using only
the theory's defining equations and allowed rules.
\end{enumerate}
An ordinary Lean theorem at level 1 does not automatically give level 3.
Even an ambient proof using no call to \texttt{classical} would not establish
that its induction is restricted to the induction scheme of $\ipra$.

\subsection{The exact repository target}

The inspected ProveIt snapshot is
\begin{center}\small\ttfamily
bb4427b10219ad61e0d2875fdde04bcd23b47ec6.
\end{center}
Its root toolchain is Lean \texttt{v4.32.0}, and its Mathlib dependency is
\texttt{v4.32.0}. The relevant modules are:
\begin{center}
\begin{tabular}{>{\raggedright\arraybackslash}p{0.31\textwidth}>{\raggedright\arraybackslash}p{0.61\textwidth}}
\toprule
Module under \texttt{ConciseIPRA} & Role verified in the source\\
\midrule
\texttt{Syntax} & Scoped terms, formulas, substitution, and the predicate
\texttt{Fm.QF}.\\
\texttt{Deduction} & Intuitionistic derivations \texttt{Prf}, including
quantifier rules and equality substitution.\\
\texttt{PR} & Function codes \texttt{PR}, terms \texttt{PTm}, the theory
\texttt{PRAAx}, and the rule \texttt{PRA.qfInd}.\\
\texttt{NativeBiInterpretation} & The equivalence
\texttt{pra\_iff\_native}, which translates a supplied internal proof.\\
\texttt{TeqRec} & A derived primitive recursor in the iterator language,
useful if implementing directly in the scalar core.\\
\bottomrule
\end{tabular}
\end{center}
The snapshot and the inspected blob hashes are recorded in
\texttt{evidence/repository\_snapshot.json}. No result about FLT has been
assumed merely because these foundational modules are present.

The file \texttt{lean/ObjectLanguageTargets.lean} defines explicit codes
for addition, multiplication, and exponentiation using
\texttt{PR.comp} and \texttt{PR.prec}. The convention is
\[
 \operatorname{addCode}(n,x)=n+x,\quad
 \operatorname{mulCode}(n,x)=nx,\quad
 \operatorname{powCode}(n,x)=x^n;
\]
the recursion argument is always first, as required by the repository.
It then builds the actual formula
\[
 \operatorname{fixedSentence}(d)=
 \forall x\,\forall y\,\forall z\,
 \bigl(x\ne0\to y\ne0\to z\ne0\to
            x^d+y^d=z^d\to\bot\bigr)
\]
using \texttt{Fm.eq}, \texttt{Fm.imp}, \texttt{Fm.bot}, and
\texttt{Fm.all}. Its definition
\[
 \texttt{InternalFLT d}
 :=\texttt{Prf PRAAx [] (fixedSentence d)}
\]
is the precise type of the missing formal proof object. It is not a
synonym for an ambient Lean proposition about natural numbers.

The uniform formula uses a free exponent, a nonzero guard, and a
disjunction of the four formulas $\exists k\,dk=n$. A nonzero multiple of
one of $3,4,5,7$ is automatically at least $3$, so this is equivalent to
the stated main family. The bound for selecting $k$ is $k\leq n$.
The code also gives the exact conditional translation to
\texttt{NPrf [] (NFm.ofPRA sentence)} using
\texttt{pra\_iff\_native}. This translation assumes an actual internal
derivation as its input; it supplies none of the arithmetic work.

\subsection{Required numerical definitions}

The file \texttt{lean/ArithmeticBlueprint.lean} records the semantic
definitions. The internal version should implement the same operations
on numerical codes, in the following order.

\paragraph{Natural arithmetic.}
Define truncated subtraction, the zero test, comparison, quotient and
remainder with a prescribed value at divisor zero, bounded sums and
products, and a first-success bounded selector. For a predicate
$\chi(i,\vec a)$, the selector returns the least $i\leq B$ with
$\chi(i,\vec a)=1$, together with a success flag. Its invariant, after
scanning $0,\ldots,j$, states that the returned candidate is either the
least success in that interval or the flag is zero and all those tests
are zero. The latter bounded universal is encoded by a product. The
invariant is therefore quantifier free after the definitions are expanded.

\paragraph{Signed arithmetic and tuples.}
Use the signed code and pairing function from
Section~\ref{sec:foundations}; implement addition, multiplication, negation,
absolute value, and sign by finite case splits. Prove that the explicit
operations satisfy the ring identities. A code for an integer is not a
quantifier over a new infinite sort: every statement expands to a numerical
formula. Rational quadratic coordinates used in the quintic proof can be
stored as integer numerators with positive integer denominators, and all
comparisons can be made after cross multiplication.

\paragraph{Divisibility and roots.}
Define the gcd computation with its clock, its Bezout accumulator, prime
testing, valuation at nonzero input, finite factor lists, and an exact-root
selector. For exponent $e\geq1$ the root selector only scans $0,\ldots,N$.
Its correctness condition is a finite equality test. Its success in the
proof follows from the coprime-factorization lemma, not from a numerical
experiment. On the few invalid inputs where a valuation is not defined
mathematically, the code returns a specified default, and every theorem
has an explicit nonzero guard.

\paragraph{The two quadratic pair rings.}
Define multiplication, conjugation, norm, rational quotient coordinates,
nearest-integer rounding, and a clocked gcd for each pair ring. Prove norm
multiplicativity by expanding polynomial expressions. The Eisenstein norm
is positive definite, so a norm bound directly bounds coordinates. The
golden norm is indefinite, so one must first implement the bounded unit
normalization from Section~\ref{sec:five}. Omitting this difference would
invalidate the bounded-factorization claim for the quintic case.

\subsection{State predicates and the central correctness contracts}

Here are the complete mathematical contents of the main state predicates.
They should be Boolean functions on encoded tuples; the corresponding
Lean predicates in the blueprint are readable specifications.

\begin{center}
\small
\begin{tabular}{>{\raggedright\arraybackslash}p{0.14\textwidth}>{\raggedright\arraybackslash}p{0.65\textwidth}>{\raggedright\arraybackslash}p{0.11\textwidth}}
\toprule
State & Predicate, in addition to the displayed equation & Rank\\
\midrule
Cubic $(x,y,z)$ & Nonzero signed pairwise coprime entries;
$x^3+y^3+z^3=0$; $3\mid z$, $3\nmid xy$. & $|z|$\\
Quartic $(x,y,z)$ & Positive entries; $\gcd(x,y)=1$;
$x^4+y^4=z^2$. & $z$\\
Quintic odd $(h,t,s,w)$ & $h\geq0$; $t,s,w>0$;
$\gcd(t,s)=1$; $t,s$ odd; $5\mid s$, $5\nmid t$;
$5^h sF(t,s)=16w^5$. & $s$\\
Quintic even $(h,t,s,w)$ & Same positivity and coprimality;
$t$ odd, $s$ even; $5\mid s$, $5\nmid t$;
$2\cdot5^h sF(t,s)=w^5$. & $s$\\
Septimal $(a,P,Q,R)$ & $a\geq1$; positive odd pairwise coprime
$P,Q,R$; equation \eqref{eq:seven-quartic}. & $a$\\
\bottomrule
\end{tabular}
\end{center}

For each of the first four rows, construct a total selector $D$ and prove
\[
 \chi(s)=1\quad\Longrightarrow\quad
       \chi(D(s))=1\ \wedge\ r(D(s))+1\leq r(s).
\tag{C}
\]
The object-level proof of (C) consists exactly of the corresponding
descent section. Invalid states return a default and do not enter its
hypotheses. The final contradiction is the shared quantifier-free
invariant (I) from Lemma~\ref{lem:clocked}.

For the septimal row one may also use the simpler fixed-box Boolean
\[
 \operatorname{bad7}(a,M)=
 \max_{1\leq P,Q,R\leq M}\chi_{D_a}(P,Q,R).
\]
Here ``max'' means an explicitly nested finite Boolean fold. Prove
$\operatorname{bad7}(1,M)=0$ and
$\operatorname{bad7}(a+1,M)=0$ from
$\operatorname{bad7}(a,M)=0$ for $a\geq1$. The fact that each new
coordinate is at most the old $R$ keeps the same $M$ throughout.
The induction formula is the literal equation
$\operatorname{bad7}(a+1,M)=0$.

In the natural-number implementation of the quartic polynomial, avoid
truncated subtraction. Use the equivalent equality
\[
 P^2+3\cdot7^4\,2^{2a}Q^2R^2
     =Q^4+7^7\,2^{4a+4}R^4.
\]
Intermediate square completions can use the signed-number code. This
small representation decision prevents the formal statement from
silently becoming a different equation.

\subsection{Construction of internal proofs, in dependency order}

The following is a finite implementation plan with identified proof
obligations. It does not call for any new number-theoretic theorem beyond
the arguments in this paper.

\begin{enumerate}
\item \textbf{Arithmetic equations.} Build internal proofs for addition,
multiplication, powers, order, and signed operations. Prove the identities
used by polynomial normalization. At this stage a normalizer may compute
the two expanded polynomials, but it must also emit a derivation using
these laws and \texttt{Prf.eqSubst}. Merely running the ambient Lean
\texttt{ring} tactic does not emit an $\ipra$ proof automatically.

\item \textbf{Bounded-search reflection.} Prove internally that the Boolean
folds and first-success selector satisfy their specified equations and
produce witnesses within their bounds. This is repeated use of
\texttt{PRA.precZero}, \texttt{PRA.precSucc}, and \texttt{PRA.qfInd}.
The invariants described above are explicit recipes for the required
\texttt{Fm.QF} proofs.

\item \textbf{Integer arithmetic tools.} Formalize
Lemma~\ref{lem:integer-tools} and primitive reduction. Use bounded lists
of prime factors; replace any induction over ``all smaller integers''
by a Boolean invariant over the finite list or bounded range. Prove the
valuation rule for a sum of unequal valuations by extracting the smaller
power and reducing the remaining equality modulo the prime.

\item \textbf{Quadratic arithmetic.} Formalize the coordinate identities,
Euclidean gcd, primality of irreducibles, finite factorization, and units.
For $E$ prove the two exact $\lambda$ divisibility tests and the six-unit
table. For $R$ prove bounded unit normalization, the classification
$\pm\varphi^j$, and the modulo-$5$ fifth-power scalar test. These are
the precise algebraic inputs to the root-extraction routines.

\item \textbf{Local transitions.} Prove (C) for cubic, quartic, and the two
quintic states. For exponent $7$, prove the base residue contradiction,
the complete four-way factor allocation, the row-III contradiction, and
the transition to parameter $a-1$. All finite allocations are to be
proved exhaustive before any branch is discharged.

\item \textbf{Finite termination.} Apply the single generic clocked-descent
derivation, or the septimal fixed-box induction. This yields the four
impossibility lemmas without an unbounded minimality predicate.

\item \textbf{Initial transformations.} Translate an arbitrary alleged
Fermat solution to a valid state. Use the explicit gcd reduction,
modulo-$9$ or modulo-$25$ calculations, two quintic parity branches,
and the septimal symmetric-polynomial and valuation argument. All
divisions must carry their already proved divisibility witnesses.

\item \textbf{Closure and export.} Construct \texttt{InternalFLT 3},
\texttt{InternalFLT 4}, \texttt{InternalFLT 5}, and \texttt{InternalFLT 7}.
Formalize exponent lifting and the four-way divisibility split to obtain
\texttt{InternalCoverage}. Finally apply \texttt{pra\_iff\_native} if a
proof in ConciseIPRA's native scalar language is desired.
\end{enumerate}

The source file declares a structure listing the required proof objects,
but does not supply an inhabitant. There are no axioms or admitted holes
pretending to fill them. The ambient blueprint similarly makes any use
of unimplemented base proofs an explicit hypothesis. This allows a future
implementation to replace each interface with a checked definition while
preserving an honest inventory of what remains.

\subsection{Build and acceptance criteria}

The \texttt{lean/README.txt} gives the commands for checking the blueprint
against the pinned ProveIt environment. A completed implementation should
be accepted only after all of the following have been checked:
\begin{enumerate}
\item The code compiles at the pinned toolchain, and the four base results
are proved without placeholders or extra axioms.
\item The completed deliverable contains the actual proof objects
\texttt{InternalFLT d}, not only the semantic statements.
\item Every inductive invariant used in those internal proofs has a
checked \texttt{Fm.QF} derivation. All algorithms used as internal
function symbols have actual \texttt{PR} codes.
\item The proof audit lists dependencies of the final declarations.
The internal proof rules must still be exactly \texttt{PRAAx} and
\texttt{Prf}; a powerful ambient fact must not enter by an added
object-language axiom.
\end{enumerate}
These criteria describe how to finish the requested formalization. They
are not represented as checks already completed by this package.

\section{Artifact verification and reproducibility}
\label{sec:verification}

The Python files serve three limited and useful roles: checking algebraic
transcription, checking finite residue tables, and enumerating the exact
exponent family. They are not proof checkers for $\ipra$ and do not replace
the human proof.

\begin{center}
\begin{tabular}{>{\raggedright\arraybackslash}p{0.36\textwidth}>{\raggedright\arraybackslash}p{0.56\textwidth}}
\toprule
Artifact & Verified role\\
\midrule
\texttt{verify\_cubic.py} & Pair-ring identities, 28,561 multiplication
and Euclidean-division samples, and the six-unit coefficient table.\\
\texttt{verify\_quintic.py} & Exact symbolic norm and fifth-power
identities; finite modulo-$2$, modulo-$5$, and modulo-$25$ tables;
bounded Euclidean and unit-normalization samples.\\
\texttt{verify\_seven.py} & Seven exact symbolic identities, including
both correction square completions and the Fermat-to-quartic substitution.\\
\texttt{coverage.py} & Exact $276/420$ residue count, exponent-reduction
witnesses, and a CSV through exponent $200$.\\
\texttt{ObjectLanguageTargets.lean} & Actual function codes and target
formulas for the internal formalization; uncompiled blueprint.\\
\texttt{ArithmeticBlueprint.lean} & Readable arithmetic definitions and
formalization interfaces; uncompiled blueprint.\\
\bottomrule
\end{tabular}
\end{center}

The algebra scripts use exact integers or rational arithmetic; the
symbolic checks use SymPy rather than floating point. Their output is
recorded in \texttt{evidence/verification.txt}. Mathematical review also
checked the unbounded arguments that sampling cannot establish: finite
factorization, the golden-ring unit normalization, all four septimal
factor allocations, the valuation extraction, the strict descent
inequalities, and the quantifier-free induction formulas.

The archive includes this article's complete modular LaTeX source,
the compiled PDF, the Lean blueprint, the verification scripts, the
source-snapshot manifest, and reproducibility instructions. The PDF was
compiled and rendered for layout inspection. The Lean files were not
compiled; there was no Lean executable in the delivery environment.
\texttt{SHA256SUMS.txt} identifies the exact files in the package.



\clearpage
\begin{thebibliography}{9}

\bibitem{proveit}
Vladimir Reshetnikov and contributors,
\emph{ProveIt: Concise presentations of intuitionistic primitive recursive
arithmetic}, snapshot
\texttt{bb4427b10219ad61e0d2875fdde04bcd23b47ec6}.
Inspected modules include \texttt{Syntax.lean}, \texttt{Deduction.lean},
\texttt{PR.lean}, \texttt{TeqRec.lean}, and
\texttt{NativeBiInterpretation.lean}.
\url{https://github.com/VladimirReshetnikov/ProveIt/tree/bb4427b10219ad61e0d2875fdde04bcd23b47ec6/Logic/PrimitiveRecursiveArithmetic/ConciseIPRA/Lean/ConciseIPRA}.

\bibitem{mathlib3}
The Mathlib community,
\emph{Mathlib.NumberTheory.FLT.Three}, documentation and source for
\texttt{fermatLastTheoremThree}.
\url{https://leanprover-community.github.io/mathlib4_docs/Mathlib/NumberTheory/FLT/Three.html}.
Consulted 15 November 2026; used as a semantic comparison, not as an
internal $\ipra$ derivation.

\bibitem{mathlib4}
The Mathlib community,
\emph{Mathlib.NumberTheory.FLT.Four}, documentation and source for
\texttt{not\_fermat\_42} and \texttt{fermatLastTheoremFour}.
\url{https://leanprover-community.github.io/mathlib4_docs/Mathlib/NumberTheory/FLT/Four.html}.
Consulted 15 November 2026.

\bibitem{five-source}
Contributors to \texttt{isomorphic-ai/fermat},
\emph{Fermat/Five}, including the modules \texttt{Coordinates},
\texttt{PowerExtraction}, \texttt{Initial}, and \texttt{Descent}.
Snapshot \texttt{07b35d10984632caa14b1af26ec5dec2363ee920}.
\url{https://github.com/isomorphic-ai/fermat/tree/07b35d10984632caa14b1af26ec5dec2363ee920/Fermat/Five}.
Inspected as a comparison source; not compiled in this package.

\bibitem{Lebesgue1840}
V.-A. Lebesgue,
\emph{D\'emonstration de l'impossibilit\'e de r\'esoudre l'\'equation
$x^7+y^7+z^7=0$ en nombres entiers},
Journal de math\'ematiques pures et appliqu\'ees, first series,
vol.~5 (1840), pp.~276--279.
\url{https://www.numdam.org/item/JMPA_1840_1_5__276_0.pdf}.

\bibitem{Lebesgue1840Add}
V.-A. Lebesgue,
\emph{Addition \`a la note sur l'\'equation $x^7+y^7+z^7=0$},
Journal de math\'ematiques pures et appliqu\'ees, first series,
vol.~5 (1840), pp.~348--349.
\url{https://www.numdam.org/item/JMPA_1840_1_5__348_0.pdf}.
This correction supplies a factor-allocation case missing from the
original proof.

\end{thebibliography}
\end{document}
